\documentclass[11pt]{article}

    \usepackage[breakable]{tcolorbox}
    \usepackage{parskip} % Stop auto-indenting (to mimic markdown behaviour)
    

    % Basic figure setup, for now with no caption control since it's done
    % automatically by Pandoc (which extracts ![](path) syntax from Markdown).
    \usepackage{graphicx}
    % Keep aspect ratio if custom image width or height is specified
    \setkeys{Gin}{keepaspectratio}
    % Maintain compatibility with old templates. Remove in nbconvert 6.0
    \let\Oldincludegraphics\includegraphics
    % Ensure that by default, figures have no caption (until we provide a
    % proper Figure object with a Caption API and a way to capture that
    % in the conversion process - todo).
    \usepackage{caption}
    \DeclareCaptionFormat{nocaption}{}
    \captionsetup{format=nocaption,aboveskip=0pt,belowskip=0pt}

    \usepackage{float}
    \floatplacement{figure}{H} % forces figures to be placed at the correct location
    \usepackage{xcolor} % Allow colors to be defined
    \usepackage{enumerate} % Needed for markdown enumerations to work
    \usepackage{geometry} % Used to adjust the document margins
    \usepackage{amsmath} % Equations
    \usepackage{amssymb} % Equations
    \usepackage{textcomp} % defines textquotesingle
    % Hack from http://tex.stackexchange.com/a/47451/13684:
    \AtBeginDocument{%
        \def\PYZsq{\textquotesingle}% Upright quotes in Pygmentized code
    }
    \usepackage{upquote} % Upright quotes for verbatim code
    \usepackage{eurosym} % defines \euro

    \usepackage{iftex}
    \ifPDFTeX
        \usepackage[T1]{fontenc}
        \IfFileExists{alphabeta.sty}{
              \usepackage{alphabeta}
          }{
              \usepackage[mathletters]{ucs}
              \usepackage[utf8x]{inputenc}
          }
    \else
        \usepackage{fontspec}
        \usepackage{unicode-math}
    \fi

    \usepackage{fancyvrb} % verbatim replacement that allows latex
    \usepackage{grffile} % extends the file name processing of package graphics
                         % to support a larger range
    \makeatletter % fix for old versions of grffile with XeLaTeX
    \@ifpackagelater{grffile}{2019/11/01}
    {
      % Do nothing on new versions
    }
    {
      \def\Gread@@xetex#1{%
        \IfFileExists{"\Gin@base".bb}%
        {\Gread@eps{\Gin@base.bb}}%
        {\Gread@@xetex@aux#1}%
      }
    }
    \makeatother
    \usepackage[Export]{adjustbox} % Used to constrain images to a maximum size
    \adjustboxset{max size={0.9\linewidth}{0.9\paperheight}}

    % The hyperref package gives us a pdf with properly built
    % internal navigation ('pdf bookmarks' for the table of contents,
    % internal cross-reference links, web links for URLs, etc.)
    \usepackage{hyperref}
    % The default LaTeX title has an obnoxious amount of whitespace. By default,
    % titling removes some of it. It also provides customization options.
    \usepackage{titling}
    \usepackage{longtable} % longtable support required by pandoc >1.10
    \usepackage{booktabs}  % table support for pandoc > 1.12.2
    \usepackage{array}     % table support for pandoc >= 2.11.3
    \usepackage{calc}      % table minipage width calculation for pandoc >= 2.11.1
    \usepackage[inline]{enumitem} % IRkernel/repr support (it uses the enumerate* environment)
    \usepackage[normalem]{ulem} % ulem is needed to support strikethroughs (\sout)
                                % normalem makes italics be italics, not underlines
    \usepackage{soul}      % strikethrough (\st) support for pandoc >= 3.0.0
    \usepackage{mathrsfs}
    

    
    % Colors for the hyperref package
    \definecolor{urlcolor}{rgb}{0,.145,.698}
    \definecolor{linkcolor}{rgb}{.71,0.21,0.01}
    \definecolor{citecolor}{rgb}{.12,.54,.11}

    % ANSI colors
    \definecolor{ansi-black}{HTML}{3E424D}
    \definecolor{ansi-black-intense}{HTML}{282C36}
    \definecolor{ansi-red}{HTML}{E75C58}
    \definecolor{ansi-red-intense}{HTML}{B22B31}
    \definecolor{ansi-green}{HTML}{00A250}
    \definecolor{ansi-green-intense}{HTML}{007427}
    \definecolor{ansi-yellow}{HTML}{DDB62B}
    \definecolor{ansi-yellow-intense}{HTML}{B27D12}
    \definecolor{ansi-blue}{HTML}{208FFB}
    \definecolor{ansi-blue-intense}{HTML}{0065CA}
    \definecolor{ansi-magenta}{HTML}{D160C4}
    \definecolor{ansi-magenta-intense}{HTML}{A03196}
    \definecolor{ansi-cyan}{HTML}{60C6C8}
    \definecolor{ansi-cyan-intense}{HTML}{258F8F}
    \definecolor{ansi-white}{HTML}{C5C1B4}
    \definecolor{ansi-white-intense}{HTML}{A1A6B2}
    \definecolor{ansi-default-inverse-fg}{HTML}{FFFFFF}
    \definecolor{ansi-default-inverse-bg}{HTML}{000000}

    % common color for the border for error outputs.
    \definecolor{outerrorbackground}{HTML}{FFDFDF}

    % commands and environments needed by pandoc snippets
    % extracted from the output of `pandoc -s`
    \providecommand{\tightlist}{%
      \setlength{\itemsep}{0pt}\setlength{\parskip}{0pt}}
    \DefineVerbatimEnvironment{Highlighting}{Verbatim}{commandchars=\\\{\}}
    % Add ',fontsize=\small' for more characters per line
    \newenvironment{Shaded}{}{}
    \newcommand{\KeywordTok}[1]{\textcolor[rgb]{0.00,0.44,0.13}{\textbf{{#1}}}}
    \newcommand{\DataTypeTok}[1]{\textcolor[rgb]{0.56,0.13,0.00}{{#1}}}
    \newcommand{\DecValTok}[1]{\textcolor[rgb]{0.25,0.63,0.44}{{#1}}}
    \newcommand{\BaseNTok}[1]{\textcolor[rgb]{0.25,0.63,0.44}{{#1}}}
    \newcommand{\FloatTok}[1]{\textcolor[rgb]{0.25,0.63,0.44}{{#1}}}
    \newcommand{\CharTok}[1]{\textcolor[rgb]{0.25,0.44,0.63}{{#1}}}
    \newcommand{\StringTok}[1]{\textcolor[rgb]{0.25,0.44,0.63}{{#1}}}
    \newcommand{\CommentTok}[1]{\textcolor[rgb]{0.38,0.63,0.69}{\textit{{#1}}}}
    \newcommand{\OtherTok}[1]{\textcolor[rgb]{0.00,0.44,0.13}{{#1}}}
    \newcommand{\AlertTok}[1]{\textcolor[rgb]{1.00,0.00,0.00}{\textbf{{#1}}}}
    \newcommand{\FunctionTok}[1]{\textcolor[rgb]{0.02,0.16,0.49}{{#1}}}
    \newcommand{\RegionMarkerTok}[1]{{#1}}
    \newcommand{\ErrorTok}[1]{\textcolor[rgb]{1.00,0.00,0.00}{\textbf{{#1}}}}
    \newcommand{\NormalTok}[1]{{#1}}

    % Additional commands for more recent versions of Pandoc
    \newcommand{\ConstantTok}[1]{\textcolor[rgb]{0.53,0.00,0.00}{{#1}}}
    \newcommand{\SpecialCharTok}[1]{\textcolor[rgb]{0.25,0.44,0.63}{{#1}}}
    \newcommand{\VerbatimStringTok}[1]{\textcolor[rgb]{0.25,0.44,0.63}{{#1}}}
    \newcommand{\SpecialStringTok}[1]{\textcolor[rgb]{0.73,0.40,0.53}{{#1}}}
    \newcommand{\ImportTok}[1]{{#1}}
    \newcommand{\DocumentationTok}[1]{\textcolor[rgb]{0.73,0.13,0.13}{\textit{{#1}}}}
    \newcommand{\AnnotationTok}[1]{\textcolor[rgb]{0.38,0.63,0.69}{\textbf{\textit{{#1}}}}}
    \newcommand{\CommentVarTok}[1]{\textcolor[rgb]{0.38,0.63,0.69}{\textbf{\textit{{#1}}}}}
    \newcommand{\VariableTok}[1]{\textcolor[rgb]{0.10,0.09,0.49}{{#1}}}
    \newcommand{\ControlFlowTok}[1]{\textcolor[rgb]{0.00,0.44,0.13}{\textbf{{#1}}}}
    \newcommand{\OperatorTok}[1]{\textcolor[rgb]{0.40,0.40,0.40}{{#1}}}
    \newcommand{\BuiltInTok}[1]{{#1}}
    \newcommand{\ExtensionTok}[1]{{#1}}
    \newcommand{\PreprocessorTok}[1]{\textcolor[rgb]{0.74,0.48,0.00}{{#1}}}
    \newcommand{\AttributeTok}[1]{\textcolor[rgb]{0.49,0.56,0.16}{{#1}}}
    \newcommand{\InformationTok}[1]{\textcolor[rgb]{0.38,0.63,0.69}{\textbf{\textit{{#1}}}}}
    \newcommand{\WarningTok}[1]{\textcolor[rgb]{0.38,0.63,0.69}{\textbf{\textit{{#1}}}}}
    \makeatletter
    \newsavebox\pandoc@box
    \newcommand*\pandocbounded[1]{%
      \sbox\pandoc@box{#1}%
      % scaling factors for width and height
      \Gscale@div\@tempa\textheight{\dimexpr\ht\pandoc@box+\dp\pandoc@box\relax}%
      \Gscale@div\@tempb\linewidth{\wd\pandoc@box}%
      % select the smaller of both
      \ifdim\@tempb\p@<\@tempa\p@
        \let\@tempa\@tempb
      \fi
      % scaling accordingly (\@tempa < 1)
      \ifdim\@tempa\p@<\p@
        \scalebox{\@tempa}{\usebox\pandoc@box}%
      % scaling not needed, use as it is
      \else
        \usebox{\pandoc@box}%
      \fi
    }
    \makeatother

    % Define a nice break command that doesn't care if a line doesn't already
    % exist.
    \def\br{\hspace*{\fill} \\* }
    % Math Jax compatibility definitions
    \def\gt{>}
    \def\lt{<}
    \let\Oldtex\TeX
    \let\Oldlatex\LaTeX
    \renewcommand{\TeX}{\textrm{\Oldtex}}
    \renewcommand{\LaTeX}{\textrm{\Oldlatex}}
    % Document parameters
    % Document title
    \title{Constellations\_memory\_demo}
    
    
    
    
    
    
    
% Pygments definitions
\makeatletter
\def\PY@reset{\let\PY@it=\relax \let\PY@bf=\relax%
    \let\PY@ul=\relax \let\PY@tc=\relax%
    \let\PY@bc=\relax \let\PY@ff=\relax}
\def\PY@tok#1{\csname PY@tok@#1\endcsname}
\def\PY@toks#1+{\ifx\relax#1\empty\else%
    \PY@tok{#1}\expandafter\PY@toks\fi}
\def\PY@do#1{\PY@bc{\PY@tc{\PY@ul{%
    \PY@it{\PY@bf{\PY@ff{#1}}}}}}}
\def\PY#1#2{\PY@reset\PY@toks#1+\relax+\PY@do{#2}}

\@namedef{PY@tok@w}{\def\PY@tc##1{\textcolor[rgb]{0.73,0.73,0.73}{##1}}}
\@namedef{PY@tok@c}{\let\PY@it=\textit\def\PY@tc##1{\textcolor[rgb]{0.24,0.48,0.48}{##1}}}
\@namedef{PY@tok@cp}{\def\PY@tc##1{\textcolor[rgb]{0.61,0.40,0.00}{##1}}}
\@namedef{PY@tok@k}{\let\PY@bf=\textbf\def\PY@tc##1{\textcolor[rgb]{0.00,0.50,0.00}{##1}}}
\@namedef{PY@tok@kp}{\def\PY@tc##1{\textcolor[rgb]{0.00,0.50,0.00}{##1}}}
\@namedef{PY@tok@kt}{\def\PY@tc##1{\textcolor[rgb]{0.69,0.00,0.25}{##1}}}
\@namedef{PY@tok@o}{\def\PY@tc##1{\textcolor[rgb]{0.40,0.40,0.40}{##1}}}
\@namedef{PY@tok@ow}{\let\PY@bf=\textbf\def\PY@tc##1{\textcolor[rgb]{0.67,0.13,1.00}{##1}}}
\@namedef{PY@tok@nb}{\def\PY@tc##1{\textcolor[rgb]{0.00,0.50,0.00}{##1}}}
\@namedef{PY@tok@nf}{\def\PY@tc##1{\textcolor[rgb]{0.00,0.00,1.00}{##1}}}
\@namedef{PY@tok@nc}{\let\PY@bf=\textbf\def\PY@tc##1{\textcolor[rgb]{0.00,0.00,1.00}{##1}}}
\@namedef{PY@tok@nn}{\let\PY@bf=\textbf\def\PY@tc##1{\textcolor[rgb]{0.00,0.00,1.00}{##1}}}
\@namedef{PY@tok@ne}{\let\PY@bf=\textbf\def\PY@tc##1{\textcolor[rgb]{0.80,0.25,0.22}{##1}}}
\@namedef{PY@tok@nv}{\def\PY@tc##1{\textcolor[rgb]{0.10,0.09,0.49}{##1}}}
\@namedef{PY@tok@no}{\def\PY@tc##1{\textcolor[rgb]{0.53,0.00,0.00}{##1}}}
\@namedef{PY@tok@nl}{\def\PY@tc##1{\textcolor[rgb]{0.46,0.46,0.00}{##1}}}
\@namedef{PY@tok@ni}{\let\PY@bf=\textbf\def\PY@tc##1{\textcolor[rgb]{0.44,0.44,0.44}{##1}}}
\@namedef{PY@tok@na}{\def\PY@tc##1{\textcolor[rgb]{0.41,0.47,0.13}{##1}}}
\@namedef{PY@tok@nt}{\let\PY@bf=\textbf\def\PY@tc##1{\textcolor[rgb]{0.00,0.50,0.00}{##1}}}
\@namedef{PY@tok@nd}{\def\PY@tc##1{\textcolor[rgb]{0.67,0.13,1.00}{##1}}}
\@namedef{PY@tok@s}{\def\PY@tc##1{\textcolor[rgb]{0.73,0.13,0.13}{##1}}}
\@namedef{PY@tok@sd}{\let\PY@it=\textit\def\PY@tc##1{\textcolor[rgb]{0.73,0.13,0.13}{##1}}}
\@namedef{PY@tok@si}{\let\PY@bf=\textbf\def\PY@tc##1{\textcolor[rgb]{0.64,0.35,0.47}{##1}}}
\@namedef{PY@tok@se}{\let\PY@bf=\textbf\def\PY@tc##1{\textcolor[rgb]{0.67,0.36,0.12}{##1}}}
\@namedef{PY@tok@sr}{\def\PY@tc##1{\textcolor[rgb]{0.64,0.35,0.47}{##1}}}
\@namedef{PY@tok@ss}{\def\PY@tc##1{\textcolor[rgb]{0.10,0.09,0.49}{##1}}}
\@namedef{PY@tok@sx}{\def\PY@tc##1{\textcolor[rgb]{0.00,0.50,0.00}{##1}}}
\@namedef{PY@tok@m}{\def\PY@tc##1{\textcolor[rgb]{0.40,0.40,0.40}{##1}}}
\@namedef{PY@tok@gh}{\let\PY@bf=\textbf\def\PY@tc##1{\textcolor[rgb]{0.00,0.00,0.50}{##1}}}
\@namedef{PY@tok@gu}{\let\PY@bf=\textbf\def\PY@tc##1{\textcolor[rgb]{0.50,0.00,0.50}{##1}}}
\@namedef{PY@tok@gd}{\def\PY@tc##1{\textcolor[rgb]{0.63,0.00,0.00}{##1}}}
\@namedef{PY@tok@gi}{\def\PY@tc##1{\textcolor[rgb]{0.00,0.52,0.00}{##1}}}
\@namedef{PY@tok@gr}{\def\PY@tc##1{\textcolor[rgb]{0.89,0.00,0.00}{##1}}}
\@namedef{PY@tok@ge}{\let\PY@it=\textit}
\@namedef{PY@tok@gs}{\let\PY@bf=\textbf}
\@namedef{PY@tok@ges}{\let\PY@bf=\textbf\let\PY@it=\textit}
\@namedef{PY@tok@gp}{\let\PY@bf=\textbf\def\PY@tc##1{\textcolor[rgb]{0.00,0.00,0.50}{##1}}}
\@namedef{PY@tok@go}{\def\PY@tc##1{\textcolor[rgb]{0.44,0.44,0.44}{##1}}}
\@namedef{PY@tok@gt}{\def\PY@tc##1{\textcolor[rgb]{0.00,0.27,0.87}{##1}}}
\@namedef{PY@tok@err}{\def\PY@bc##1{{\setlength{\fboxsep}{\string -\fboxrule}\fcolorbox[rgb]{1.00,0.00,0.00}{1,1,1}{\strut ##1}}}}
\@namedef{PY@tok@kc}{\let\PY@bf=\textbf\def\PY@tc##1{\textcolor[rgb]{0.00,0.50,0.00}{##1}}}
\@namedef{PY@tok@kd}{\let\PY@bf=\textbf\def\PY@tc##1{\textcolor[rgb]{0.00,0.50,0.00}{##1}}}
\@namedef{PY@tok@kn}{\let\PY@bf=\textbf\def\PY@tc##1{\textcolor[rgb]{0.00,0.50,0.00}{##1}}}
\@namedef{PY@tok@kr}{\let\PY@bf=\textbf\def\PY@tc##1{\textcolor[rgb]{0.00,0.50,0.00}{##1}}}
\@namedef{PY@tok@bp}{\def\PY@tc##1{\textcolor[rgb]{0.00,0.50,0.00}{##1}}}
\@namedef{PY@tok@fm}{\def\PY@tc##1{\textcolor[rgb]{0.00,0.00,1.00}{##1}}}
\@namedef{PY@tok@vc}{\def\PY@tc##1{\textcolor[rgb]{0.10,0.09,0.49}{##1}}}
\@namedef{PY@tok@vg}{\def\PY@tc##1{\textcolor[rgb]{0.10,0.09,0.49}{##1}}}
\@namedef{PY@tok@vi}{\def\PY@tc##1{\textcolor[rgb]{0.10,0.09,0.49}{##1}}}
\@namedef{PY@tok@vm}{\def\PY@tc##1{\textcolor[rgb]{0.10,0.09,0.49}{##1}}}
\@namedef{PY@tok@sa}{\def\PY@tc##1{\textcolor[rgb]{0.73,0.13,0.13}{##1}}}
\@namedef{PY@tok@sb}{\def\PY@tc##1{\textcolor[rgb]{0.73,0.13,0.13}{##1}}}
\@namedef{PY@tok@sc}{\def\PY@tc##1{\textcolor[rgb]{0.73,0.13,0.13}{##1}}}
\@namedef{PY@tok@dl}{\def\PY@tc##1{\textcolor[rgb]{0.73,0.13,0.13}{##1}}}
\@namedef{PY@tok@s2}{\def\PY@tc##1{\textcolor[rgb]{0.73,0.13,0.13}{##1}}}
\@namedef{PY@tok@sh}{\def\PY@tc##1{\textcolor[rgb]{0.73,0.13,0.13}{##1}}}
\@namedef{PY@tok@s1}{\def\PY@tc##1{\textcolor[rgb]{0.73,0.13,0.13}{##1}}}
\@namedef{PY@tok@mb}{\def\PY@tc##1{\textcolor[rgb]{0.40,0.40,0.40}{##1}}}
\@namedef{PY@tok@mf}{\def\PY@tc##1{\textcolor[rgb]{0.40,0.40,0.40}{##1}}}
\@namedef{PY@tok@mh}{\def\PY@tc##1{\textcolor[rgb]{0.40,0.40,0.40}{##1}}}
\@namedef{PY@tok@mi}{\def\PY@tc##1{\textcolor[rgb]{0.40,0.40,0.40}{##1}}}
\@namedef{PY@tok@il}{\def\PY@tc##1{\textcolor[rgb]{0.40,0.40,0.40}{##1}}}
\@namedef{PY@tok@mo}{\def\PY@tc##1{\textcolor[rgb]{0.40,0.40,0.40}{##1}}}
\@namedef{PY@tok@ch}{\let\PY@it=\textit\def\PY@tc##1{\textcolor[rgb]{0.24,0.48,0.48}{##1}}}
\@namedef{PY@tok@cm}{\let\PY@it=\textit\def\PY@tc##1{\textcolor[rgb]{0.24,0.48,0.48}{##1}}}
\@namedef{PY@tok@cpf}{\let\PY@it=\textit\def\PY@tc##1{\textcolor[rgb]{0.24,0.48,0.48}{##1}}}
\@namedef{PY@tok@c1}{\let\PY@it=\textit\def\PY@tc##1{\textcolor[rgb]{0.24,0.48,0.48}{##1}}}
\@namedef{PY@tok@cs}{\let\PY@it=\textit\def\PY@tc##1{\textcolor[rgb]{0.24,0.48,0.48}{##1}}}

\def\PYZbs{\char`\\}
\def\PYZus{\char`\_}
\def\PYZob{\char`\{}
\def\PYZcb{\char`\}}
\def\PYZca{\char`\^}
\def\PYZam{\char`\&}
\def\PYZlt{\char`\<}
\def\PYZgt{\char`\>}
\def\PYZsh{\char`\#}
\def\PYZpc{\char`\%}
\def\PYZdl{\char`\$}
\def\PYZhy{\char`\-}
\def\PYZsq{\char`\'}
\def\PYZdq{\char`\"}
\def\PYZti{\char`\~}
% for compatibility with earlier versions
\def\PYZat{@}
\def\PYZlb{[}
\def\PYZrb{]}
\makeatother


    % For linebreaks inside Verbatim environment from package fancyvrb.
    \makeatletter
        \newbox\Wrappedcontinuationbox
        \newbox\Wrappedvisiblespacebox
        \newcommand*\Wrappedvisiblespace {\textcolor{red}{\textvisiblespace}}
        \newcommand*\Wrappedcontinuationsymbol {\textcolor{red}{\llap{\tiny$\m@th\hookrightarrow$}}}
        \newcommand*\Wrappedcontinuationindent {3ex }
        \newcommand*\Wrappedafterbreak {\kern\Wrappedcontinuationindent\copy\Wrappedcontinuationbox}
        % Take advantage of the already applied Pygments mark-up to insert
        % potential linebreaks for TeX processing.
        %        {, <, #, %, $, ' and ": go to next line.
        %        _, }, ^, &, >, - and ~: stay at end of broken line.
        % Use of \textquotesingle for straight quote.
        \newcommand*\Wrappedbreaksatspecials {%
            \def\PYGZus{\discretionary{\char`\_}{\Wrappedafterbreak}{\char`\_}}%
            \def\PYGZob{\discretionary{}{\Wrappedafterbreak\char`\{}{\char`\{}}%
            \def\PYGZcb{\discretionary{\char`\}}{\Wrappedafterbreak}{\char`\}}}%
            \def\PYGZca{\discretionary{\char`\^}{\Wrappedafterbreak}{\char`\^}}%
            \def\PYGZam{\discretionary{\char`\&}{\Wrappedafterbreak}{\char`\&}}%
            \def\PYGZlt{\discretionary{}{\Wrappedafterbreak\char`\<}{\char`\<}}%
            \def\PYGZgt{\discretionary{\char`\>}{\Wrappedafterbreak}{\char`\>}}%
            \def\PYGZsh{\discretionary{}{\Wrappedafterbreak\char`\#}{\char`\#}}%
            \def\PYGZpc{\discretionary{}{\Wrappedafterbreak\char`\%}{\char`\%}}%
            \def\PYGZdl{\discretionary{}{\Wrappedafterbreak\char`\$}{\char`\$}}%
            \def\PYGZhy{\discretionary{\char`\-}{\Wrappedafterbreak}{\char`\-}}%
            \def\PYGZsq{\discretionary{}{\Wrappedafterbreak\textquotesingle}{\textquotesingle}}%
            \def\PYGZdq{\discretionary{}{\Wrappedafterbreak\char`\"}{\char`\"}}%
            \def\PYGZti{\discretionary{\char`\~}{\Wrappedafterbreak}{\char`\~}}%
        }
        % Some characters . , ; ? ! / are not pygmentized.
        % This macro makes them "active" and they will insert potential linebreaks
        \newcommand*\Wrappedbreaksatpunct {%
            \lccode`\~`\.\lowercase{\def~}{\discretionary{\hbox{\char`\.}}{\Wrappedafterbreak}{\hbox{\char`\.}}}%
            \lccode`\~`\,\lowercase{\def~}{\discretionary{\hbox{\char`\,}}{\Wrappedafterbreak}{\hbox{\char`\,}}}%
            \lccode`\~`\;\lowercase{\def~}{\discretionary{\hbox{\char`\;}}{\Wrappedafterbreak}{\hbox{\char`\;}}}%
            \lccode`\~`\:\lowercase{\def~}{\discretionary{\hbox{\char`\:}}{\Wrappedafterbreak}{\hbox{\char`\:}}}%
            \lccode`\~`\?\lowercase{\def~}{\discretionary{\hbox{\char`\?}}{\Wrappedafterbreak}{\hbox{\char`\?}}}%
            \lccode`\~`\!\lowercase{\def~}{\discretionary{\hbox{\char`\!}}{\Wrappedafterbreak}{\hbox{\char`\!}}}%
            \lccode`\~`\/\lowercase{\def~}{\discretionary{\hbox{\char`\/}}{\Wrappedafterbreak}{\hbox{\char`\/}}}%
            \catcode`\.\active
            \catcode`\,\active
            \catcode`\;\active
            \catcode`\:\active
            \catcode`\?\active
            \catcode`\!\active
            \catcode`\/\active
            \lccode`\~`\~
        }
    \makeatother

    \let\OriginalVerbatim=\Verbatim
    \makeatletter
    \renewcommand{\Verbatim}[1][1]{%
        %\parskip\z@skip
        \sbox\Wrappedcontinuationbox {\Wrappedcontinuationsymbol}%
        \sbox\Wrappedvisiblespacebox {\FV@SetupFont\Wrappedvisiblespace}%
        \def\FancyVerbFormatLine ##1{\hsize\linewidth
            \vtop{\raggedright\hyphenpenalty\z@\exhyphenpenalty\z@
                \doublehyphendemerits\z@\finalhyphendemerits\z@
                \strut ##1\strut}%
        }%
        % If the linebreak is at a space, the latter will be displayed as visible
        % space at end of first line, and a continuation symbol starts next line.
        % Stretch/shrink are however usually zero for typewriter font.
        \def\FV@Space {%
            \nobreak\hskip\z@ plus\fontdimen3\font minus\fontdimen4\font
            \discretionary{\copy\Wrappedvisiblespacebox}{\Wrappedafterbreak}
            {\kern\fontdimen2\font}%
        }%

        % Allow breaks at special characters using \PYG... macros.
        \Wrappedbreaksatspecials
        % Breaks at punctuation characters . , ; ? ! and / need catcode=\active
        \OriginalVerbatim[#1,codes*=\Wrappedbreaksatpunct]%
    }
    \makeatother

    % Exact colors from NB
    \definecolor{incolor}{HTML}{303F9F}
    \definecolor{outcolor}{HTML}{D84315}
    \definecolor{cellborder}{HTML}{CFCFCF}
    \definecolor{cellbackground}{HTML}{F7F7F7}

    % prompt
    \makeatletter
    \newcommand{\boxspacing}{\kern\kvtcb@left@rule\kern\kvtcb@boxsep}
    \makeatother
    \newcommand{\prompt}[4]{
        {\ttfamily\llap{{\color{#2}[#3]:\hspace{3pt}#4}}\vspace{-\baselineskip}}
    }
    

    
    % Prevent overflowing lines due to hard-to-break entities
    \sloppy
    % Setup hyperref package
    \hypersetup{
      breaklinks=true,  % so long urls are correctly broken across lines
      colorlinks=true,
      urlcolor=urlcolor,
      linkcolor=linkcolor,
      citecolor=citecolor,
      }
    % Slightly bigger margins than the latex defaults
    
    \geometry{verbose,tmargin=1in,bmargin=1in,lmargin=1in,rmargin=1in}
    
    

\begin{document}
    
    \maketitle
    
    

    
    \hypertarget{harey-constellation-cards}{%
\section{HARey constellation cards}\label{harey-constellation-cards}}

\textbf{HARey} (pen name of Hans Augusto Reyersbach) was an illustrator
and author. He took an interest in astronomy and wrote the books
\emph{The Stars: A New Way to see them} and \emph{Find the
constellations} to teach children about stargazing. In them he tried to
redrawn the constellations diagrams to be more similar to the objects
they are named after, instead of the abstract shapes used in modern
astronomy.

This project is my personal homage to HARey's work, and it borrows his
idea of making learning the constellations a challenge by creating a
Memory-like game to recognize the constellations without the lines.

    \begin{tcolorbox}[breakable, size=fbox, boxrule=1pt, pad at break*=1mm,colback=cellbackground, colframe=cellborder]
\prompt{In}{incolor}{26}{\boxspacing}
\begin{Verbatim}[commandchars=\\\{\}]
\PY{k+kn}{import} \PY{n+nn}{numpy} \PY{k}{as} \PY{n+nn}{np}
\PY{k+kn}{import} \PY{n+nn}{matplotlib}\PY{n+nn}{.}\PY{n+nn}{pyplot} \PY{k}{as} \PY{n+nn}{plt}
\PY{k+kn}{from} \PY{n+nn}{matplotlib}\PY{n+nn}{.}\PY{n+nn}{font\PYZus{}manager} \PY{k+kn}{import} \PY{n}{FontProperties}
\PY{k+kn}{from} \PY{n+nn}{datetime} \PY{k+kn}{import} \PY{n}{datetime}
\PY{k+kn}{import} \PY{n+nn}{pytz}
\PY{k+kn}{import} \PY{n+nn}{pandas} \PY{k}{as} \PY{n+nn}{pd}
\end{Verbatim}
\end{tcolorbox}

    \hypertarget{the-project-base}{%
\subsection{The project base}\label{the-project-base}}

This project borrows a lot from \textbf{Stellarium}, an open source
planetary software. In the \emph{Sky Culture} section users can see
different constellations shapes created by different cultures, and among
them there are HARey's drawings. The \emph{index.json} file provides the
constellation lines data and asterisms and the Hipparcos catalogues
provides the stars positions to draw the constellations.

    The \textbf{constellations} dictionary contains each constellations
lines and stars. The key is an abbreviation of the constellation name
(e.g., Ori for Orion). Constellations with keys like \textbf{.Orib} are
a subgroup of stars that are assigned a name by HARey (e.g.~.Orib is
Orion's belt). The \textbf{main\_ids} list contains the ids of the 88
IAU constellations.

    \begin{tcolorbox}[breakable, size=fbox, boxrule=1pt, pad at break*=1mm,colback=cellbackground, colframe=cellborder]
\prompt{In}{incolor}{27}{\boxspacing}
\begin{Verbatim}[commandchars=\\\{\}]
\PY{c+c1}{\PYZsh{}Load the constellations  }
\PY{n}{constellations}\PY{p}{,} \PY{n}{main\PYZus{}ids}\PY{p}{,} \PY{n}{asterisms}\PY{p}{,} \PY{n}{helpers}\PY{p}{,} \PY{n}{named\PYZus{}stars} \PY{o}{=} \PY{n}{load\PYZus{}constellations}\PY{p}{(}\PY{p}{)}
\PY{c+c1}{\PYZsh{} Load the stars positions and the constellations they belong to}
\PY{n}{stars} \PY{o}{=} \PY{n}{load\PYZus{}stars}\PY{p}{(}\PY{p}{)}

\PY{c+c1}{\PYZsh{} Print the first 10 ids as an example}
\PY{n+nb}{print}\PY{p}{(}\PY{n}{main\PYZus{}ids}\PY{p}{[}\PY{l+m+mi}{0}\PY{p}{:}\PY{l+m+mi}{10}\PY{p}{]}\PY{p}{)}
\end{Verbatim}
\end{tcolorbox}

    \begin{Verbatim}[commandchars=\\\{\}]
['UMi', 'Dra', 'Cas', 'Cep', 'Cam', 'UMa', 'Leo', 'CVn', 'LMi', 'Boo']
    \end{Verbatim}

    The stars are projected using a \textbf{stereographic} projection, which
is a projection from spherical surface to plane which conserves the
angles at which lines meet, locally preserving the shapes (further from
the center, the stretching distorts the overall diagram). The center of
the projection is the brightest star of each constellation.

    \begin{tcolorbox}[breakable, size=fbox, boxrule=1pt, pad at break*=1mm,colback=cellbackground, colframe=cellborder]
\prompt{In}{incolor}{28}{\boxspacing}
\begin{Verbatim}[commandchars=\\\{\}]
\PY{n}{stars} \PY{o}{=} \PY{n}{stars}\PY{p}{[}\PY{n}{stars}\PY{p}{[}\PY{l+s+s1}{\PYZsq{}}\PY{l+s+s1}{magnitude}\PY{l+s+s1}{\PYZsq{}}\PY{p}{]} \PY{o}{\PYZlt{}} \PY{l+m+mi}{5}\PY{p}{]}
\PY{n}{local\PYZus{}stars} \PY{o}{=} \PY{n}{stars}\PY{o}{.}\PY{n}{loc}\PY{p}{[}\PY{n}{constellations}\PY{p}{[}\PY{l+s+s1}{\PYZsq{}}\PY{l+s+s1}{UMa}\PY{l+s+s1}{\PYZsq{}}\PY{p}{]}\PY{p}{[}\PY{l+s+s1}{\PYZsq{}}\PY{l+s+s1}{stars}\PY{l+s+s1}{\PYZsq{}}\PY{p}{]}\PY{p}{]}

\PY{c+c1}{\PYZsh{} Take the brightest star of the constellation as the ccenter (zenith) of the projection}
\PY{n}{brightest\PYZus{}star} \PY{o}{=} \PY{n}{local\PYZus{}stars}\PY{o}{.}\PY{n}{iloc}\PY{p}{[}\PY{n}{np}\PY{o}{.}\PY{n}{argmin}\PY{p}{(}\PY{n}{local\PYZus{}stars}\PY{p}{[}\PY{l+s+s1}{\PYZsq{}}\PY{l+s+s1}{magnitude}\PY{l+s+s1}{\PYZsq{}}\PY{p}{]}\PY{p}{)}\PY{p}{]}

\PY{c+c1}{\PYZsh{} Perform the stereographic projection}
\PY{n}{stars\PYZus{}x}\PY{p}{,} \PY{n}{stars\PYZus{}y} \PY{o}{=} \PY{n}{stereo\PYZus{}centered}\PY{p}{(}\PY{n}{stars}\PY{p}{[}\PY{l+s+s1}{\PYZsq{}}\PY{l+s+s1}{ra}\PY{l+s+s1}{\PYZsq{}}\PY{p}{]}\PY{p}{,} \PY{n}{stars}\PY{p}{[}\PY{l+s+s1}{\PYZsq{}}\PY{l+s+s1}{dec}\PY{l+s+s1}{\PYZsq{}}\PY{p}{]}\PY{p}{,} \PY{n}{brightest\PYZus{}star}\PY{p}{[}\PY{l+s+s1}{\PYZsq{}}\PY{l+s+s1}{ra}\PY{l+s+s1}{\PYZsq{}}\PY{p}{]}\PY{p}{,} \PY{n}{brightest\PYZus{}star}\PY{p}{[}\PY{l+s+s1}{\PYZsq{}}\PY{l+s+s1}{dec}\PY{l+s+s1}{\PYZsq{}}\PY{p}{]}\PY{p}{)}

\PY{c+c1}{\PYZsh{} Convert the values to  Pandas series by adding the index}
\PY{n}{stars\PYZus{}x} \PY{o}{=} \PY{n}{pd}\PY{o}{.}\PY{n}{Series}\PY{p}{(}\PY{n}{data} \PY{o}{=} \PY{n}{stars\PYZus{}x}\PY{p}{,} \PY{n}{index}\PY{o}{=}\PY{n}{stars}\PY{o}{.}\PY{n}{index}\PY{p}{)}
\PY{n}{stars\PYZus{}y} \PY{o}{=} \PY{n}{pd}\PY{o}{.}\PY{n}{Series}\PY{p}{(}\PY{n}{data} \PY{o}{=} \PY{n}{stars\PYZus{}y}\PY{p}{,} \PY{n}{index}\PY{o}{=}\PY{n}{stars}\PY{o}{.}\PY{n}{index}\PY{p}{)}
\end{Verbatim}
\end{tcolorbox}

    \begin{tcolorbox}[breakable, size=fbox, boxrule=1pt, pad at break*=1mm,colback=cellbackground, colframe=cellborder]
\prompt{In}{incolor}{29}{\boxspacing}
\begin{Verbatim}[commandchars=\\\{\}]
\PY{n}{fig}\PY{p}{,} \PY{p}{(}\PY{n}{ax1}\PY{p}{,} \PY{n}{ax2}\PY{p}{)} \PY{o}{=} \PY{n}{plt}\PY{o}{.}\PY{n}{subplots}\PY{p}{(}\PY{n}{facecolor}\PY{o}{=}\PY{l+s+s1}{\PYZsq{}}\PY{l+s+s1}{midnightblue}\PY{l+s+s1}{\PYZsq{}}\PY{p}{,} \PY{n}{ncols} \PY{o}{=} \PY{l+m+mi}{2}\PY{p}{,} \PY{n}{figsize}\PY{o}{=}\PY{p}{(}\PY{l+m+mi}{10}\PY{p}{,}\PY{l+m+mi}{5}\PY{p}{)}\PY{p}{)}
\PY{n}{fig}\PY{o}{.}\PY{n}{suptitle}\PY{p}{(}\PY{l+s+s1}{\PYZsq{}}\PY{l+s+s1}{Ursa Major}\PY{l+s+s1}{\PYZsq{}}\PY{p}{,} \PY{n}{color}\PY{o}{=}\PY{l+s+s1}{\PYZsq{}}\PY{l+s+s1}{w}\PY{l+s+s1}{\PYZsq{}}\PY{p}{,} \PY{n}{fontsize} \PY{o}{=} \PY{l+m+mi}{20}\PY{p}{)}
\PY{n}{ax1}\PY{o}{.}\PY{n}{set\PYZus{}aspect}\PY{p}{(}\PY{l+s+s1}{\PYZsq{}}\PY{l+s+s1}{equal}\PY{l+s+s1}{\PYZsq{}}\PY{p}{)}
\PY{n}{ax1}\PY{o}{.}\PY{n}{set\PYZus{}axis\PYZus{}off}\PY{p}{(}\PY{p}{)}
\PY{n}{ax1}\PY{o}{.}\PY{n}{set\PYZus{}xlim}\PY{p}{(}\PY{o}{\PYZhy{}}\PY{l+m+mf}{0.1}\PY{p}{,}\PY{l+m+mf}{0.3}\PY{p}{)}
\PY{n}{ax1}\PY{o}{.}\PY{n}{set\PYZus{}ylim}\PY{p}{(}\PY{o}{\PYZhy{}}\PY{l+m+mf}{0.2}\PY{p}{,}\PY{l+m+mf}{0.2}\PY{p}{)}

\PY{c+c1}{\PYZsh{} The stars size is proportional to its brightness}
\PY{n}{ax1}\PY{o}{.}\PY{n}{scatter}\PY{p}{(}\PY{n}{stars\PYZus{}x}\PY{p}{,} \PY{n}{stars\PYZus{}y}\PY{p}{,} \PY{n}{s}\PY{o}{=}\PY{l+m+mi}{50}\PY{o}{*}\PY{l+m+mi}{10}\PY{o}{*}\PY{o}{*}\PY{p}{(}\PY{n}{stars}\PY{p}{[}\PY{l+s+s1}{\PYZsq{}}\PY{l+s+s1}{magnitude}\PY{l+s+s1}{\PYZsq{}}\PY{p}{]}\PY{o}{/}\PY{o}{\PYZhy{}}\PY{l+m+mf}{2.5}\PY{p}{)}\PY{p}{,} \PY{n}{color}\PY{o}{=}\PY{l+s+s1}{\PYZsq{}}\PY{l+s+s1}{w}\PY{l+s+s1}{\PYZsq{}}\PY{p}{)}

\PY{n}{ax2}\PY{o}{.}\PY{n}{set\PYZus{}aspect}\PY{p}{(}\PY{l+s+s1}{\PYZsq{}}\PY{l+s+s1}{equal}\PY{l+s+s1}{\PYZsq{}}\PY{p}{)}
\PY{n}{ax2}\PY{o}{.}\PY{n}{set\PYZus{}axis\PYZus{}off}\PY{p}{(}\PY{p}{)}
\PY{n}{ax2}\PY{o}{.}\PY{n}{set\PYZus{}xlim}\PY{p}{(}\PY{o}{\PYZhy{}}\PY{l+m+mf}{0.1}\PY{p}{,}\PY{l+m+mf}{0.3}\PY{p}{)}
\PY{n}{ax2}\PY{o}{.}\PY{n}{set\PYZus{}ylim}\PY{p}{(}\PY{o}{\PYZhy{}}\PY{l+m+mf}{0.2}\PY{p}{,}\PY{l+m+mf}{0.2}\PY{p}{)}

\PY{c+c1}{\PYZsh{} Plot the lines of the constellations Ursa Major}
\PY{k}{for} \PY{n}{line} \PY{o+ow}{in} \PY{n}{constellations}\PY{p}{[}\PY{l+s+s1}{\PYZsq{}}\PY{l+s+s1}{UMa}\PY{l+s+s1}{\PYZsq{}}\PY{p}{]}\PY{p}{[}\PY{l+s+s1}{\PYZsq{}}\PY{l+s+s1}{lines}\PY{l+s+s1}{\PYZsq{}}\PY{p}{]}\PY{p}{:}
    \PY{n}{ax2}\PY{o}{.}\PY{n}{plot}\PY{p}{(}\PY{n}{stars\PYZus{}x}\PY{p}{[}\PY{n}{line}\PY{p}{]}\PY{p}{,} \PY{n}{stars\PYZus{}y}\PY{p}{[}\PY{n}{line}\PY{p}{]}\PY{p}{,} \PY{n}{color}\PY{o}{=}\PY{l+s+s1}{\PYZsq{}}\PY{l+s+s1}{y}\PY{l+s+s1}{\PYZsq{}}\PY{p}{,} \PY{n}{linewidth}\PY{o}{=}\PY{l+m+mf}{0.5}\PY{p}{)}
    
\PY{n}{ax2}\PY{o}{.}\PY{n}{scatter}\PY{p}{(}\PY{n}{stars\PYZus{}x}\PY{p}{,} \PY{n}{stars\PYZus{}y}\PY{p}{,} \PY{n}{s}\PY{o}{=}\PY{l+m+mi}{50}\PY{o}{*}\PY{l+m+mi}{10}\PY{o}{*}\PY{o}{*}\PY{p}{(}\PY{n}{stars}\PY{p}{[}\PY{l+s+s1}{\PYZsq{}}\PY{l+s+s1}{magnitude}\PY{l+s+s1}{\PYZsq{}}\PY{p}{]}\PY{o}{/}\PY{o}{\PYZhy{}}\PY{l+m+mf}{2.5}\PY{p}{)}\PY{p}{,} \PY{n}{color}\PY{o}{=}\PY{l+s+s1}{\PYZsq{}}\PY{l+s+s1}{w}\PY{l+s+s1}{\PYZsq{}}\PY{p}{)}
\PY{n}{plt}\PY{o}{.}\PY{n}{show}\PY{p}{(}\PY{p}{)}
\end{Verbatim}
\end{tcolorbox}

    \begin{center}
    \adjustimage{max size={0.9\linewidth}{0.9\paperheight}}{Constellations_memory_demo_files/Constellations_memory_demo_7_0.png}
    \end{center}
    { \hspace*{\fill} \\}
    
    \hypertarget{the-hareymain-module}{%
\section{The HAReyMain module}\label{the-hareymain-module}}

The process of creating the constellations plots is automated by the
\textbf{HAReyMain} module, which creates plots similar to HARey's
drawings.

To make cards that can be used by everyone in their native language, I
added the \textbf{names.csv} file, which contains the keys to identify
the objects inside the HARey module and the translated names. The
default language is IAU-EN, which is the standard IAU denomination of
objects; ENGLISH uses the names that HARey wrote in his book. For now
only ENGLISH and ITALIAN are supported, but I'd like some help to add
more languages.

The names also include newlines to wrap around long ones (this works in
the Matplotlib plots but not in InkScape, see the SIS\_SCRIPT section).

    \begin{tcolorbox}[breakable, size=fbox, boxrule=1pt, pad at break*=1mm,colback=cellbackground, colframe=cellborder]
\prompt{In}{incolor}{30}{\boxspacing}
\begin{Verbatim}[commandchars=\\\{\}]
\PY{k+kn}{from} \PY{n+nn}{HARey}\PY{n+nn}{.}\PY{n+nn}{harey\PYZus{}main} \PY{k+kn}{import} \PY{n}{HAReyMain}
\end{Verbatim}
\end{tcolorbox}

    \begin{tcolorbox}[breakable, size=fbox, boxrule=1pt, pad at break*=1mm,colback=cellbackground, colframe=cellborder]
\prompt{In}{incolor}{31}{\boxspacing}
\begin{Verbatim}[commandchars=\\\{\}]
\PY{c+c1}{\PYZsh{} Read the constellations file, the HIP catalogue and the object names}
\PY{n}{harey} \PY{o}{=} \PY{n}{HAReyMain}\PY{p}{(}\PY{n}{language}\PY{o}{=}\PY{l+s+s1}{\PYZsq{}}\PY{l+s+s1}{ENGLISH}\PY{l+s+s1}{\PYZsq{}}\PY{p}{,} \PY{n}{star\PYZus{}colors}\PY{o}{=}\PY{l+s+s1}{\PYZsq{}}\PY{l+s+s1}{helland}\PY{l+s+s1}{\PYZsq{}}\PY{p}{)}

\PY{c+c1}{\PYZsh{} Set two custom fonts for the plots and the cards names}
\PY{n}{labels\PYZus{}font} \PY{o}{=} \PY{n}{FontProperties}\PY{p}{(}\PY{n}{fname}\PY{o}{=}\PY{l+s+s1}{\PYZsq{}}\PY{l+s+s1}{fonts/PatrickHand\PYZhy{}Regular.ttf}\PY{l+s+s1}{\PYZsq{}}\PY{p}{)}
\PY{n}{cardback\PYZus{}font} \PY{o}{=} \PY{n}{FontProperties}\PY{p}{(}\PY{n}{fname}\PY{o}{=}\PY{l+s+s1}{\PYZsq{}}\PY{l+s+s1}{fonts/TheWalkyr.ttf}\PY{l+s+s1}{\PYZsq{}}\PY{p}{)}

\PY{n}{harey}\PY{o}{.}\PY{n}{set\PYZus{}fonts}\PY{p}{(}\PY{p}{\PYZob{}}\PY{l+s+s1}{\PYZsq{}}\PY{l+s+s1}{labels}\PY{l+s+s1}{\PYZsq{}}\PY{p}{:}\PY{n}{labels\PYZus{}font}\PY{p}{,} \PY{l+s+s1}{\PYZsq{}}\PY{l+s+s1}{cardback}\PY{l+s+s1}{\PYZsq{}}\PY{p}{:}\PY{n}{cardback\PYZus{}font}\PY{p}{\PYZcb{}}\PY{p}{)}
\end{Verbatim}
\end{tcolorbox}

    \begin{Verbatim}[commandchars=\\\{\}]
Loading constellations diagrams{\ldots}     Done!
Loading star coordinates{\ldots}     Done!
Computing stars colors{\ldots}     Done!
Loading custom markers{\ldots}       Done!
Loading the object names{\ldots}       Done!


Using the tarot-round format, 2.75x4.75 in, using the template at
/home/menegattig/Desktop/HARey/HARey/cardbacks/tarot\_round.png
    \end{Verbatim}

    \hypertarget{star-sizes-and-colors}{%
\subsection{Star sizes and colors}\label{star-sizes-and-colors}}

In his books, HARey drew different markers for different magnitudes,
making it easier to see the magnitude of each star, whereas modern maps
all use round stars.

The size of a star depends on its brightness, which is computed from its
magnitude \(M\) as \(B=B_0\cdot 10 ^{-0.4M}\).

Different functions are used by many softwares depending on the
objective. Most set the size equal to the brightness, but digging a lot
into the math I found that the visual size of a star is given by the
Airy disk, in which case it is \(\propto (M-M_l)\), where \(M_l\), the
limit magnitude, is the faintest magnitude visible.

The python library Skyfield uses \(\propto (M-M_l)^2\), and I finally
settled for \(\propto (M-M_l)^{1.5}\), which I felt was a good
compromise considering that the different symbols already help
distinguishing the bright stars.

The faintest stars in the Hipparcos catalogue are of magnitude 13, but
the naked eye cannot see beyond 6-6.5, and depending on the light
pollution the value could be much lower. The default value used by the
harey library is 6.0, to prioritize a clear plot instead of a fancy one.
It can be changed wit \texttt{harey.set\_limiting\_magnitude()}

    \begin{tcolorbox}[breakable, size=fbox, boxrule=1pt, pad at break*=1mm,colback=cellbackground, colframe=cellborder]
\prompt{In}{incolor}{32}{\boxspacing}
\begin{Verbatim}[commandchars=\\\{\}]
\PY{n}{fig} \PY{o}{=} \PY{n}{harey}\PY{o}{.}\PY{n}{plot\PYZus{}legend}\PY{p}{(}\PY{p}{)}
\PY{n}{plt}\PY{o}{.}\PY{n}{show}\PY{p}{(}\PY{p}{)}
\end{Verbatim}
\end{tcolorbox}

    \begin{center}
    \adjustimage{max size={0.9\linewidth}{0.9\paperheight}}{Constellations_memory_demo_files/Constellations_memory_demo_12_0.png}
    \end{center}
    { \hspace*{\fill} \\}
    
    For a bit of flare, each star has its own color, defined by the B-V
color index. This value is converted in a RGB color using a colormap
present in the Stellarium software.

    \begin{tcolorbox}[breakable, size=fbox, boxrule=1pt, pad at break*=1mm,colback=cellbackground, colframe=cellborder]
\prompt{In}{incolor}{33}{\boxspacing}
\begin{Verbatim}[commandchars=\\\{\}]
\PY{n}{fig} \PY{o}{=} \PY{n}{harey}\PY{o}{.}\PY{n}{plot\PYZus{}star\PYZus{}cmap}\PY{p}{(}\PY{p}{)}
\PY{n}{plt}\PY{o}{.}\PY{n}{show}\PY{p}{(}\PY{p}{)}
\end{Verbatim}
\end{tcolorbox}

    \begin{center}
    \adjustimage{max size={0.9\linewidth}{0.9\paperheight}}{Constellations_memory_demo_files/Constellations_memory_demo_14_0.png}
    \end{center}
    { \hspace*{\fill} \\}
    
    The colors used in the plots are shared by all the plotting functions
and can be changed with
\texttt{harey.set\_colors(\{\textquotesingle{}name\textquotesingle{}:color\})}.

    \begin{tcolorbox}[breakable, size=fbox, boxrule=1pt, pad at break*=1mm,colback=cellbackground, colframe=cellborder]
\prompt{In}{incolor}{34}{\boxspacing}
\begin{Verbatim}[commandchars=\\\{\}]
\PY{n+nb}{print}\PY{p}{(}\PY{n}{harey}\PY{o}{.}\PY{n}{colors}\PY{p}{)}
\end{Verbatim}
\end{tcolorbox}

    \begin{Verbatim}[commandchars=\\\{\}]
\{'star': 'white', 'constellations': 'white', 'sky': 'xkcd:midnight', 'ecliptic':
'crimson', 'horizon': 'white', 'cardinal\_markers': 'darkred', 'grid': 'yellow',
'asterisms': 'limegreen', 'helpers': 'coral', 'starmap\_border': 'xkcd:gold',
'star\_labels': 'gold', 'constellation\_labels': 'cyan', 'ecliptic\_label':
'crimson', 'asterism\_labels': 'lime', 'constellation\_parts': 'violet',
'horizon\_label': 'white', 'mater': 'xkcd:light blue', 'cardback\_1': 'xkcd:marine
blue', 'cardback\_2': 'xkcd:blood', 'accent\_1': 'darkgoldenrod', 'accent\_2':
'darkgoldenrod'\}
    \end{Verbatim}

    \hypertarget{plotting-options}{%
\subsubsection{Plotting options}\label{plotting-options}}

The plotting functions share most of these boolean flags to enable or
disable plotting options. They are saved in the \texttt{harey.flags}
dictionary. They can be changed before a plot with
\texttt{harey.set\_flags\_\_(\{\textquotesingle{}flag\textquotesingle{}:value\})},
then they are reset to default values automatically after the plot.

\begin{itemize}
\tightlist
\item
  \textbf{CON\_LINES} (\emph{default False}) : Plot the constellation
  lines
\item
  \textbf{STAR\_COLORS} (\emph{default False}): Use the real star colors
  in the plots. Otherwise, all the stars will be white
  \textbf{CON\_NAMES} (\emph{default False}): Show the constellation
  names
\item
  \textbf{CON\_PARTS} (\emph{default False}): Show the constellation
  parts (e.g.~feet, head, body, ecc.) from the HARey diagrams
\item
  \textbf{STAR\_NAMES} (\emph{default False}): Plot the brightest stars
  names
\item
  \textbf{ASTERISMS} (\emph{default False}): Plot the asterisms lines.
  Asterisms are patterns of stars that are easy to recognize but not a
  constellation (e.g.~the Big Dipper)
\item
  \textbf{HELPERS} (\emph{default False}): Plot the helper lines. These
  are imaginary lines that connect stars to each other and make it
  easier to find features in the sky
\item
  \textbf{GRID} (\emph{default False}): Plot the RA-dec grid in
  equatorial and polar maps
\item
  \textbf{HAREY\_MAERKERS} (\emph{default True}): Use the HARey markers
  instead of points
\item
  \textbf{SIS\_SCRIPT} (\emph{default False}): Creates a Simple InkScape
  Script for the labels. See the next section.
\item
  \textbf{SHOW} (\emph{default True}): Show the image (default True)
\item
  \textbf{SAVE} (\emph{default False}): Save the figure with a default
  name. If a different \textbf{save\_name} is provided, the plot is
  saved automatically
\end{itemize}

In a similar ways, the default color used are saved in the
\textbf{harey\_colors} dictionary, and changed with
\textbf{harey.set\_colors()}.

    \hypertarget{the-sis_script-option}{%
\subsubsection{The SIS\_SCRIPT option}\label{the-sis_script-option}}

Plotting the labels proved a more difficult task than expected as there
is no simple way to put the labels without overlapping other labels or
objects. I've tried the python libraries \emph{adjustText} and
\emph{textalloc} that fix this issue by moving each label away and then
try to dynamically simulate their movement towards their initial
position while avoiding everything else, but found that it took too long
to simulate the movements and the result was not satisfactory (also,
there is no way to rotate the labels, which drastically improve
readability).

For this reason I decided to manually fix the labels before printing the
cards. To simplify this tedious task, I used InkScape with the extension
\textbf{Simple Inkscape Scripting}, which enables python-like scripts to
perform repetitive tasks in InkScape.

To do this, enable the SIS\_SCRIPT flag and a new \emph{image\_name.py}
file will be created inside the folder inkscape\_scripts containing the
labels informations. Now open the image (saved without labels) in
Inkscape and resize the canvas to the image (file → Document Properties
→ Resize to content). Then go to Extensions → Render → Simple InkScape
Scripting and select the \emph{image\_name.py} file and launch it with
Apply.

All the labels will appear in the image in the right position and color,
and fixing them will hopefully require only some resizing and a small
translation. Text on newlines will be split in two text windows.

    \hypertarget{constellations-cards}{%
\subsection{Constellations cards}\label{constellations-cards}}

The constellations are projected around the brightest star using a
stereographic projection that preserves the shapes.

When the \textbf{BEST\_AR} option is enabled, the stars are rotated to
minimize the aspect ratio (x-spread over y-spread) and fill as much as
possible the card space. Otherwise, the card is plotted North up. Stars
that are not part of the constellations are plotted with a little
transparency to make it easier to see the constellation.

The card size is set by \texttt{harey.set\_card\_template()}. Right now,
the format supported are poker-round (2.5x3.5 in), tarot-round
(2.75x4.75 in), jumbo-round (3.5x5.5 in) and their square variants.

Here is shown how to plot the constellation of Orion (ID `Ori')

    \begin{tcolorbox}[breakable, size=fbox, boxrule=1pt, pad at break*=1mm,colback=cellbackground, colframe=cellborder]
\prompt{In}{incolor}{35}{\boxspacing}
\begin{Verbatim}[commandchars=\\\{\}]
\PY{c+c1}{\PYZsh{} Using a card template different than the default one}
\PY{n}{harey}\PY{o}{.}\PY{n}{set\PYZus{}card\PYZus{}template}\PY{p}{(}\PY{l+s+s1}{\PYZsq{}}\PY{l+s+s1}{poker\PYZhy{}round}\PY{l+s+s1}{\PYZsq{}}\PY{p}{)}

\PY{n}{harey}\PY{o}{.}\PY{n}{set\PYZus{}flags}\PY{p}{(}\PY{p}{\PYZob{}}\PY{l+s+s1}{\PYZsq{}}\PY{l+s+s1}{CON\PYZus{}LINES}\PY{l+s+s1}{\PYZsq{}}\PY{p}{:}\PY{k+kc}{True}\PY{p}{\PYZcb{}}\PY{p}{)}
\PY{n}{harey}\PY{o}{.}\PY{n}{plot\PYZus{}card}\PY{p}{(}\PY{l+s+s1}{\PYZsq{}}\PY{l+s+s1}{Ori}\PY{l+s+s1}{\PYZsq{}}\PY{p}{,}\PY{n}{BEST\PYZus{}AR}\PY{o}{=}\PY{k+kc}{True}\PY{p}{)}
\end{Verbatim}
\end{tcolorbox}

    \begin{Verbatim}[commandchars=\\\{\}]
Using the poker-round format, 2.50x3.50 in, using the template at
/home/menegattig/Desktop/HARey/HARey/cardbacks/poker\_round.png
    \end{Verbatim}

    \begin{center}
    \adjustimage{max size={0.9\linewidth}{0.9\paperheight}}{Constellations_memory_demo_files/Constellations_memory_demo_20_1.png}
    \end{center}
    { \hspace*{\fill} \\}
    
    Another plot of orion, but with no lines and colored stars

    \begin{tcolorbox}[breakable, size=fbox, boxrule=1pt, pad at break*=1mm,colback=cellbackground, colframe=cellborder]
\prompt{In}{incolor}{36}{\boxspacing}
\begin{Verbatim}[commandchars=\\\{\}]
\PY{n}{harey}\PY{o}{.}\PY{n}{set\PYZus{}flags}\PY{p}{(}\PY{p}{\PYZob{}}\PY{l+s+s1}{\PYZsq{}}\PY{l+s+s1}{STAR\PYZus{}COLORS}\PY{l+s+s1}{\PYZsq{}}\PY{p}{:}\PY{k+kc}{True}\PY{p}{\PYZcb{}}\PY{p}{)}
\PY{n}{harey}\PY{o}{.}\PY{n}{plot\PYZus{}card}\PY{p}{(}\PY{l+s+s1}{\PYZsq{}}\PY{l+s+s1}{Ori}\PY{l+s+s1}{\PYZsq{}}\PY{p}{,} \PY{n}{BEST\PYZus{}AR}\PY{o}{=}\PY{k+kc}{True}\PY{p}{)}
\end{Verbatim}
\end{tcolorbox}

    \begin{center}
    \adjustimage{max size={0.9\linewidth}{0.9\paperheight}}{Constellations_memory_demo_files/Constellations_memory_demo_22_0.png}
    \end{center}
    { \hspace*{\fill} \\}
    
    orion again, but disabling the Harey markers and using simple points to
plot the constellation. Also the limiting magnitude is set to 10.

    \begin{tcolorbox}[breakable, size=fbox, boxrule=1pt, pad at break*=1mm,colback=cellbackground, colframe=cellborder]
\prompt{In}{incolor}{37}{\boxspacing}
\begin{Verbatim}[commandchars=\\\{\}]
\PY{n}{harey}\PY{o}{.}\PY{n}{set\PYZus{}limiting\PYZus{}magnitude}\PY{p}{(}\PY{l+m+mi}{10}\PY{p}{)}
\PY{n}{harey}\PY{o}{.}\PY{n}{set\PYZus{}flags}\PY{p}{(}\PY{p}{\PYZob{}}\PY{l+s+s1}{\PYZsq{}}\PY{l+s+s1}{HAREY\PYZus{}MARKERS}\PY{l+s+s1}{\PYZsq{}}\PY{p}{:}\PY{k+kc}{False}\PY{p}{,} \PY{l+s+s1}{\PYZsq{}}\PY{l+s+s1}{STAR\PYZus{}COLORS}\PY{l+s+s1}{\PYZsq{}}\PY{p}{:}\PY{k+kc}{True}\PY{p}{\PYZcb{}}\PY{p}{)}
\PY{n}{harey}\PY{o}{.}\PY{n}{plot\PYZus{}card}\PY{p}{(}\PY{l+s+s1}{\PYZsq{}}\PY{l+s+s1}{Ori}\PY{l+s+s1}{\PYZsq{}}\PY{p}{,} \PY{n}{BEST\PYZus{}AR}\PY{o}{=}\PY{k+kc}{True}\PY{p}{,} \PY{n}{star\PYZus{}size}\PY{o}{=}\PY{l+m+mi}{70}\PY{p}{)}
\end{Verbatim}
\end{tcolorbox}

    \begin{center}
    \adjustimage{max size={0.9\linewidth}{0.9\paperheight}}{Constellations_memory_demo_files/Constellations_memory_demo_24_0.png}
    \end{center}
    { \hspace*{\fill} \\}
    
    Orion again, constellation line, names of the constellation parts and
names of the stars. The labels position is set to be in the middle,
which does not look good, so I recommend to manually adjust the lables
in InkScape when printing a card set.

    \begin{tcolorbox}[breakable, size=fbox, boxrule=1pt, pad at break*=1mm,colback=cellbackground, colframe=cellborder]
\prompt{In}{incolor}{38}{\boxspacing}
\begin{Verbatim}[commandchars=\\\{\}]
\PY{n}{harey}\PY{o}{.}\PY{n}{set\PYZus{}limiting\PYZus{}magnitude}\PY{p}{(}\PY{l+m+mi}{6}\PY{p}{)}
\PY{n}{harey}\PY{o}{.}\PY{n}{set\PYZus{}flags}\PY{p}{(}\PY{p}{\PYZob{}}\PY{l+s+s1}{\PYZsq{}}\PY{l+s+s1}{CON\PYZus{}LINES}\PY{l+s+s1}{\PYZsq{}}\PY{p}{:}\PY{k+kc}{True}\PY{p}{,} \PY{l+s+s1}{\PYZsq{}}\PY{l+s+s1}{CON\PYZus{}PARTS}\PY{l+s+s1}{\PYZsq{}}\PY{p}{:}\PY{k+kc}{True}\PY{p}{,} \PY{l+s+s1}{\PYZsq{}}\PY{l+s+s1}{STAR\PYZus{}NAMES}\PY{l+s+s1}{\PYZsq{}}\PY{p}{:}\PY{k+kc}{True}\PY{p}{,} \PY{l+s+s1}{\PYZsq{}}\PY{l+s+s1}{STAR\PYZus{}COLORS}\PY{l+s+s1}{\PYZsq{}}\PY{p}{:}\PY{k+kc}{True}\PY{p}{\PYZcb{}}\PY{p}{)}
\PY{n}{harey}\PY{o}{.}\PY{n}{plot\PYZus{}card}\PY{p}{(}\PY{l+s+s1}{\PYZsq{}}\PY{l+s+s1}{Ori}\PY{l+s+s1}{\PYZsq{}}\PY{p}{,} \PY{n}{BEST\PYZus{}AR}\PY{o}{=}\PY{k+kc}{True}\PY{p}{)}
\end{Verbatim}
\end{tcolorbox}

    \begin{center}
    \adjustimage{max size={0.9\linewidth}{0.9\paperheight}}{Constellations_memory_demo_files/Constellations_memory_demo_26_0.png}
    \end{center}
    { \hspace*{\fill} \\}
    
    \texttt{plot\_cardback} creates the cardback for the constellations,
colorizing the black\_and\_white image of the cardback and adding the
name of the constellation to the image.

    \begin{tcolorbox}[breakable, size=fbox, boxrule=1pt, pad at break*=1mm,colback=cellbackground, colframe=cellborder]
\prompt{In}{incolor}{39}{\boxspacing}
\begin{Verbatim}[commandchars=\\\{\}]
\PY{n}{harey}\PY{o}{.}\PY{n}{plot\PYZus{}cardback}\PY{p}{(}\PY{l+s+s1}{\PYZsq{}}\PY{l+s+s1}{Ori}\PY{l+s+s1}{\PYZsq{}}\PY{p}{)}
\end{Verbatim}
\end{tcolorbox}

    \begin{center}
    \adjustimage{max size={0.9\linewidth}{0.9\paperheight}}{Constellations_memory_demo_files/Constellations_memory_demo_28_0.png}
    \end{center}
    { \hspace*{\fill} \\}
    
    \hypertarget{printing-the-cards}{%
\subsubsection{Printing the cards}\label{printing-the-cards}}

The actual print is done by first creating all the sets inside a folder,
the using the print-and-play method to arrange all the cards inside a
pdf. \texttt{print\_card\_set()} creates a card set for one
constellation (an image of the constellation without lines, one with
lines, and two different cardbacks).

Here SIS\_SCRIPT is enabled, so cards will be saved without names and a
new \textbf{inkscape\_scripts} folder will be created containing the
\emph{.py} scripts to create the labels in inkscape.

    \begin{tcolorbox}[breakable, size=fbox, boxrule=1pt, pad at break*=1mm,colback=cellbackground, colframe=cellborder]
\prompt{In}{incolor}{40}{\boxspacing}
\begin{Verbatim}[commandchars=\\\{\}]
\PY{n}{harey}\PY{o}{.}\PY{n}{set\PYZus{}card\PYZus{}template}\PY{p}{(}\PY{l+s+s1}{\PYZsq{}}\PY{l+s+s1}{jumbo\PYZhy{}round}\PY{l+s+s1}{\PYZsq{}}\PY{p}{)}
\PY{n}{plot\PYZus{}constellations} \PY{o}{=} \PY{p}{[}\PY{l+s+s1}{\PYZsq{}}\PY{l+s+s1}{Ori}\PY{l+s+s1}{\PYZsq{}}\PY{p}{,} \PY{l+s+s1}{\PYZsq{}}\PY{l+s+s1}{Gem}\PY{l+s+s1}{\PYZsq{}}\PY{p}{]}

\PY{k}{for} \PY{n+nb}{id} \PY{o+ow}{in} \PY{n}{plot\PYZus{}constellations}\PY{p}{:}
	\PY{n}{harey}\PY{o}{.}\PY{n}{set\PYZus{}flags}\PY{p}{(}\PY{p}{\PYZob{}}\PY{l+s+s1}{\PYZsq{}}\PY{l+s+s1}{STAR\PYZus{}COLORS}\PY{l+s+s1}{\PYZsq{}}\PY{p}{:}\PY{k+kc}{True}\PY{p}{,} \PY{l+s+s1}{\PYZsq{}}\PY{l+s+s1}{STAR\PYZus{}NAMES}\PY{l+s+s1}{\PYZsq{}}\PY{p}{:}\PY{k+kc}{True}\PY{p}{,} \PY{l+s+s1}{\PYZsq{}}\PY{l+s+s1}{CON\PYZus{}PARTS}\PY{l+s+s1}{\PYZsq{}}\PY{p}{:}\PY{k+kc}{True}\PY{p}{,} \PY{l+s+s1}{\PYZsq{}}\PY{l+s+s1}{SIS\PYZus{}SCRIPT}\PY{l+s+s1}{\PYZsq{}}\PY{p}{:}\PY{k+kc}{False}\PY{p}{\PYZcb{}}\PY{p}{)}
	\PY{n}{harey}\PY{o}{.}\PY{n}{print\PYZus{}card\PYZus{}set}\PY{p}{(}\PY{n+nb}{id}\PY{p}{,} \PY{n}{save\PYZus{}folder}\PY{o}{=}\PY{l+s+s1}{\PYZsq{}}\PY{l+s+s1}{Card\PYZus{}output}\PY{l+s+s1}{\PYZsq{}}\PY{p}{,} \PY{n}{BEST\PYZus{}AR}\PY{o}{=}\PY{k+kc}{True}\PY{p}{,} \PY{n}{bleed}\PY{o}{=}\PY{l+m+mf}{0.0}\PY{p}{)}
\end{Verbatim}
\end{tcolorbox}

    \begin{Verbatim}[commandchars=\\\{\}]
Using the jumbo-round format, 3.50x5.50 in, using the template at
/home/menegattig/Desktop/HARey/HARey/cardbacks/jumbo\_round.png
Created card set for Ori
Created card set for Gem
    \end{Verbatim}

    \texttt{print\_and\_play()} arranges the cards inside the folder in a
pdf document ready to print. The print is intended to be 2-sided, and to
avoid small misalignement between the fronts and the backs of the cards,
the formers have been saved with a bleed of 0.1 inches, aa small extra
area to completely cover the card backs, where the cutting should take
place.

The print process can drastically change the color of the cards, so
doing several trials before printing more cards is recommended.

    \begin{tcolorbox}[breakable, size=fbox, boxrule=1pt, pad at break*=1mm,colback=cellbackground, colframe=cellborder]
\prompt{In}{incolor}{41}{\boxspacing}
\begin{Verbatim}[commandchars=\\\{\}]
\PY{n}{harey}\PY{o}{.}\PY{n}{print\PYZus{}and\PYZus{}play}\PY{p}{(}\PY{n}{folder} \PY{o}{=} \PY{l+s+s1}{\PYZsq{}}\PY{l+s+s1}{Card\PYZus{}output}\PY{l+s+s1}{\PYZsq{}}\PY{p}{,} \PY{n}{filename}\PY{o}{=}\PY{l+s+s1}{\PYZsq{}}\PY{l+s+s1}{Cards.pdf}\PY{l+s+s1}{\PYZsq{}}\PY{p}{,} \PY{n}{bleed}\PY{o}{=}\PY{l+m+mf}{0.0}\PY{p}{)}
\end{Verbatim}
\end{tcolorbox}

    \begin{Verbatim}[commandchars=\\\{\}]

32 cards have been printed in the file Cards.pdf

    \end{Verbatim}

    \hypertarget{asterisms-and-helper-lines}{%
\subsection{Asterisms and Helper
lines}\label{asterisms-and-helper-lines}}

Aterisms and helper lines are plotted by the function
\texttt{harey.asterism\_plot()}. Asterism are group of stars that are
easy to recognize in the sky but are not constellation by themselves.
Helper rays, with IDs from `HR01' to `HR14', are lines that connect
bright stars and make it easier to find stars in the sky, the most
famous one of which is the one that connects the Big Dipper Guardians
stars to the North Star.

    \begin{tcolorbox}[breakable, size=fbox, boxrule=1pt, pad at break*=1mm,colback=cellbackground, colframe=cellborder]
\prompt{In}{incolor}{42}{\boxspacing}
\begin{Verbatim}[commandchars=\\\{\}]
\PY{c+c1}{\PYZsh{} Print asterism names}
\PY{n+nb}{print}\PY{p}{(}\PY{l+s+s1}{\PYZsq{}}\PY{l+s+s1}{Asterism IDs and names}\PY{l+s+s1}{\PYZsq{}}\PY{p}{)}
\PY{k}{for} \PY{n+nb}{id} \PY{o+ow}{in} \PY{n}{harey}\PY{o}{.}\PY{n}{asterisms}\PY{o}{.}\PY{n}{keys}\PY{p}{(}\PY{p}{)}\PY{p}{:}
    \PY{n+nb}{print}\PY{p}{(}\PY{l+s+sa}{f}\PY{l+s+s1}{\PYZsq{}}\PY{l+s+si}{\PYZob{}}\PY{n+nb}{id}\PY{l+s+si}{\PYZcb{}}\PY{l+s+s1}{ : }\PY{l+s+si}{\PYZob{}}\PY{n}{harey}\PY{o}{.}\PY{n}{names}\PY{p}{[}\PY{n+nb}{id}\PY{p}{]}\PY{o}{.}\PY{n}{replace}\PY{p}{(}\PY{l+s+s1}{\PYZsq{}}\PY{l+s+se}{\PYZbs{}n}\PY{l+s+s1}{\PYZsq{}}\PY{p}{,}\PY{+w}{ }\PY{l+s+s1}{\PYZsq{}}\PY{l+s+s1}{ }\PY{l+s+s1}{\PYZsq{}}\PY{p}{)}\PY{l+s+si}{\PYZcb{}}\PY{l+s+s1}{\PYZsq{}}\PY{p}{)}
\end{Verbatim}
\end{tcolorbox}

    \begin{Verbatim}[commandchars=\\\{\}]
Asterism IDs and names
BigD : Big Dipper
VirD : Virgin's Diamond
GSq : Great Square
GHx : Great Hexagon
SerCd : Serpent's tail
SerCp : Serpent's head
SumT : Summer Triangle
CruF : False Cross
    \end{Verbatim}

    \begin{tcolorbox}[breakable, size=fbox, boxrule=1pt, pad at break*=1mm,colback=cellbackground, colframe=cellborder]
\prompt{In}{incolor}{43}{\boxspacing}
\begin{Verbatim}[commandchars=\\\{\}]
\PY{n}{harey}\PY{o}{.}\PY{n}{set\PYZus{}flags}\PY{p}{(}\PY{p}{\PYZob{}}\PY{l+s+s1}{\PYZsq{}}\PY{l+s+s1}{CON\PYZus{}LINES}\PY{l+s+s1}{\PYZsq{}}\PY{p}{:}\PY{k+kc}{True}\PY{p}{,} \PY{l+s+s1}{\PYZsq{}}\PY{l+s+s1}{CON\PYZus{}NAMES}\PY{l+s+s1}{\PYZsq{}}\PY{p}{:}\PY{k+kc}{True}\PY{p}{\PYZcb{}}\PY{p}{)}
\PY{n}{harey}\PY{o}{.}\PY{n}{plot\PYZus{}asterism}\PY{p}{(}\PY{l+s+s1}{\PYZsq{}}\PY{l+s+s1}{HR01}\PY{l+s+s1}{\PYZsq{}}\PY{p}{)}
\end{Verbatim}
\end{tcolorbox}

    \begin{center}
    \adjustimage{max size={0.9\linewidth}{0.9\paperheight}}{Constellations_memory_demo_files/Constellations_memory_demo_35_0.png}
    \end{center}
    { \hspace*{\fill} \\}
    
    \begin{tcolorbox}[breakable, size=fbox, boxrule=1pt, pad at break*=1mm,colback=cellbackground, colframe=cellborder]
\prompt{In}{incolor}{44}{\boxspacing}
\begin{Verbatim}[commandchars=\\\{\}]
\PY{n}{harey}\PY{o}{.}\PY{n}{set\PYZus{}flags}\PY{p}{(}\PY{p}{\PYZob{}}\PY{l+s+s1}{\PYZsq{}}\PY{l+s+s1}{CON\PYZus{}LINES}\PY{l+s+s1}{\PYZsq{}}\PY{p}{:}\PY{k+kc}{True}\PY{p}{,} \PY{l+s+s1}{\PYZsq{}}\PY{l+s+s1}{CON\PYZus{}NAMES}\PY{l+s+s1}{\PYZsq{}}\PY{p}{:}\PY{k+kc}{True}\PY{p}{,} \PY{l+s+s1}{\PYZsq{}}\PY{l+s+s1}{ZODIAC}\PY{l+s+s1}{\PYZsq{}}\PY{p}{:}\PY{k+kc}{True}\PY{p}{\PYZcb{}}\PY{p}{)}
\PY{n}{harey}\PY{o}{.}\PY{n}{plot\PYZus{}asterism}\PY{p}{(}\PY{l+s+s1}{\PYZsq{}}\PY{l+s+s1}{VirD}\PY{l+s+s1}{\PYZsq{}}\PY{p}{)}
\end{Verbatim}
\end{tcolorbox}

    \begin{center}
    \adjustimage{max size={0.9\linewidth}{0.9\paperheight}}{Constellations_memory_demo_files/Constellations_memory_demo_36_0.png}
    \end{center}
    { \hspace*{\fill} \\}
    
    \hypertarget{sky-view}{%
\subsection{Sky View}\label{sky-view}}

The \texttt{plot\_sky\_view} method creates an alt-az view of the sky at
a given time and location, with the same plotting options as before.

In the sky views and sky maps, changes in the font and the overall size
of the plot may change dramatically the stars and labels appearance. It
may be necessary to change the \texttt{star\_size} and
\texttt{font\_sizes} options to achieve a good result.

The sky view here is plot with aterism kines (green) and helper lines
(orange)

    \begin{tcolorbox}[breakable, size=fbox, boxrule=1pt, pad at break*=1mm,colback=cellbackground, colframe=cellborder]
\prompt{In}{incolor}{45}{\boxspacing}
\begin{Verbatim}[commandchars=\\\{\}]
\PY{n}{datetime\PYZus{}str} \PY{o}{=} \PY{l+s+s1}{\PYZsq{}}\PY{l+s+s1}{27\PYZhy{}3\PYZhy{}2025 21:30}\PY{l+s+s1}{\PYZsq{}}

\PY{n}{date} \PY{o}{=} \PY{n}{datetime}\PY{o}{.}\PY{n}{strptime}\PY{p}{(}\PY{n}{datetime\PYZus{}str}\PY{p}{,} \PY{l+s+s1}{\PYZsq{}}\PY{l+s+si}{\PYZpc{}d}\PY{l+s+s1}{\PYZhy{}}\PY{l+s+s1}{\PYZpc{}}\PY{l+s+s1}{m\PYZhy{}}\PY{l+s+s1}{\PYZpc{}}\PY{l+s+s1}{Y }\PY{l+s+s1}{\PYZpc{}}\PY{l+s+s1}{H:}\PY{l+s+s1}{\PYZpc{}}\PY{l+s+s1}{M}\PY{l+s+s1}{\PYZsq{}}\PY{p}{)}
\PY{c+c1}{\PYZsh{} Observer takes as input the latitude and longitude in decimal format (not sessagesimal)}
\PY{n}{observer} \PY{o}{=} \PY{n}{harey}\PY{o}{.}\PY{n}{Observer}\PY{p}{(}\PY{l+s+s1}{\PYZsq{}}\PY{l+s+s1}{45 N}\PY{l+s+s1}{\PYZsq{}}\PY{p}{,}\PY{l+s+s1}{\PYZsq{}}\PY{l+s+s1}{11 E}\PY{l+s+s1}{\PYZsq{}}\PY{p}{)}
\PY{c+c1}{\PYZsh{} Set the observation time and timezone}
\PY{n}{observer}\PY{o}{.}\PY{n}{at\PYZus{}time}\PY{p}{(}\PY{n}{date}\PY{p}{,} \PY{n}{pytz}\PY{o}{.}\PY{n}{timezone}\PY{p}{(}\PY{l+s+s1}{\PYZsq{}}\PY{l+s+s1}{Europe/Rome}\PY{l+s+s1}{\PYZsq{}}\PY{p}{)}\PY{p}{)}

\PY{n}{harey}\PY{o}{.}\PY{n}{set\PYZus{}flags}\PY{p}{(}\PY{p}{\PYZob{}}\PY{l+s+s1}{\PYZsq{}}\PY{l+s+s1}{STAR\PYZus{}COLORS}\PY{l+s+s1}{\PYZsq{}}\PY{p}{:}\PY{k+kc}{True}\PY{p}{,} \PY{l+s+s1}{\PYZsq{}}\PY{l+s+s1}{STAR\PYZus{}NAMES}\PY{l+s+s1}{\PYZsq{}}\PY{p}{:}\PY{k+kc}{True}\PY{p}{,} \PY{l+s+s1}{\PYZsq{}}\PY{l+s+s1}{ASTERISMS}\PY{l+s+s1}{\PYZsq{}}\PY{p}{:}\PY{k+kc}{True}\PY{p}{\PYZcb{}}\PY{p}{)}
\PY{n}{harey}\PY{o}{.}\PY{n}{plot\PYZus{}sky\PYZus{}view}\PY{p}{(}\PY{n}{observer}\PY{p}{)}
\end{Verbatim}
\end{tcolorbox}

    \begin{center}
    \adjustimage{max size={0.9\linewidth}{0.9\paperheight}}{Constellations_memory_demo_files/Constellations_memory_demo_38_0.png}
    \end{center}
    { \hspace*{\fill} \\}
    
    The same sky view, but with constellation lines, names and parts.

    \begin{tcolorbox}[breakable, size=fbox, boxrule=1pt, pad at break*=1mm,colback=cellbackground, colframe=cellborder]
\prompt{In}{incolor}{46}{\boxspacing}
\begin{Verbatim}[commandchars=\\\{\}]
\PY{n}{harey}\PY{o}{.}\PY{n}{set\PYZus{}flags}\PY{p}{(}\PY{p}{\PYZob{}}\PY{l+s+s1}{\PYZsq{}}\PY{l+s+s1}{CON\PYZus{}LINES}\PY{l+s+s1}{\PYZsq{}}\PY{p}{:}\PY{k+kc}{True}\PY{p}{,} \PY{l+s+s1}{\PYZsq{}}\PY{l+s+s1}{CON\PYZus{}NAMES}\PY{l+s+s1}{\PYZsq{}}\PY{p}{:}\PY{k+kc}{True}\PY{p}{,} \PY{l+s+s1}{\PYZsq{}}\PY{l+s+s1}{STAR\PYZus{}NAMES}\PY{l+s+s1}{\PYZsq{}}\PY{p}{:}\PY{k+kc}{True}\PY{p}{,} \PY{l+s+s1}{\PYZsq{}}\PY{l+s+s1}{CON\PYZus{}PARTS}\PY{l+s+s1}{\PYZsq{}}\PY{p}{:}\PY{k+kc}{True}\PY{p}{,} \PY{l+s+s1}{\PYZsq{}}\PY{l+s+s1}{STAR\PYZus{}COLORS}\PY{l+s+s1}{\PYZsq{}}\PY{p}{:}\PY{k+kc}{True}\PY{p}{\PYZcb{}}\PY{p}{)}
\PY{n}{harey}\PY{o}{.}\PY{n}{plot\PYZus{}sky\PYZus{}view}\PY{p}{(}\PY{n}{observer}\PY{p}{,} \PY{n}{font\PYZus{}sizes}\PY{o}{=}\PY{p}{(}\PY{l+m+mi}{5}\PY{p}{,}\PY{l+m+mi}{7}\PY{p}{)}\PY{p}{)}
\end{Verbatim}
\end{tcolorbox}

    \begin{center}
    \adjustimage{max size={0.9\linewidth}{0.9\paperheight}}{Constellations_memory_demo_files/Constellations_memory_demo_40_0.png}
    \end{center}
    { \hspace*{\fill} \\}
    
    \hypertarget{polar-map-of-the-sky}{%
\subsection{Polar map of the sky}\label{polar-map-of-the-sky}}

The polar map uses either a stereographic projection, which preserves
angles but deforms areas moving away from the center, or an azimuthal
projection, which limits the enlargement of constellations but deforms
the shapes.

    \begin{tcolorbox}[breakable, size=fbox, boxrule=1pt, pad at break*=1mm,colback=cellbackground, colframe=cellborder]
\prompt{In}{incolor}{47}{\boxspacing}
\begin{Verbatim}[commandchars=\\\{\}]
\PY{n}{harey}\PY{o}{.}\PY{n}{set\PYZus{}flags}\PY{p}{(}\PY{p}{\PYZob{}}\PY{l+s+s1}{\PYZsq{}}\PY{l+s+s1}{CON\PYZus{}LINES}\PY{l+s+s1}{\PYZsq{}}\PY{p}{:}\PY{k+kc}{True}\PY{p}{,} \PY{l+s+s1}{\PYZsq{}}\PY{l+s+s1}{GRID}\PY{l+s+s1}{\PYZsq{}}\PY{p}{:}\PY{k+kc}{True}\PY{p}{,} \PY{l+s+s1}{\PYZsq{}}\PY{l+s+s1}{CON\PYZus{}NAMES}\PY{l+s+s1}{\PYZsq{}}\PY{p}{:}\PY{k+kc}{True}\PY{p}{,} \PY{l+s+s1}{\PYZsq{}}\PY{l+s+s1}{ASTERISMS}\PY{l+s+s1}{\PYZsq{}}\PY{p}{:}\PY{k+kc}{True}\PY{p}{,} \PY{l+s+s1}{\PYZsq{}}\PY{l+s+s1}{STAR\PYZus{}NAMES}\PY{l+s+s1}{\PYZsq{}}\PY{p}{:}\PY{k+kc}{True}\PY{p}{,} \PY{l+s+s1}{\PYZsq{}}\PY{l+s+s1}{SIS\PYZus{}SCRIPT}\PY{l+s+s1}{\PYZsq{}}\PY{p}{:}\PY{k+kc}{True}\PY{p}{,} \PY{l+s+s1}{\PYZsq{}}\PY{l+s+s1}{ADD CALENDAR}\PY{l+s+s1}{\PYZsq{}}\PY{p}{:}\PY{k+kc}{True}\PY{p}{\PYZcb{}}\PY{p}{)}
\PY{n}{harey}\PY{o}{.}\PY{n}{polar\PYZus{}map}\PY{p}{(}\PY{n}{pole} \PY{o}{=} \PY{l+s+s1}{\PYZsq{}}\PY{l+s+s1}{N}\PY{l+s+s1}{\PYZsq{}}\PY{p}{)}
\end{Verbatim}
\end{tcolorbox}

    \begin{center}
    \adjustimage{max size={0.9\linewidth}{0.9\paperheight}}{Constellations_memory_demo_files/Constellations_memory_demo_42_0.png}
    \end{center}
    { \hspace*{\fill} \\}
    
    A polar map of the southern polar sky, using the azimuthal projection.

    \begin{tcolorbox}[breakable, size=fbox, boxrule=1pt, pad at break*=1mm,colback=cellbackground, colframe=cellborder]
\prompt{In}{incolor}{48}{\boxspacing}
\begin{Verbatim}[commandchars=\\\{\}]
\PY{n}{harey}\PY{o}{.}\PY{n}{set\PYZus{}flags}\PY{p}{(}\PY{p}{\PYZob{}}\PY{l+s+s1}{\PYZsq{}}\PY{l+s+s1}{CON\PYZus{}LINES}\PY{l+s+s1}{\PYZsq{}}\PY{p}{:}\PY{k+kc}{True}\PY{p}{,} \PY{l+s+s1}{\PYZsq{}}\PY{l+s+s1}{GRID}\PY{l+s+s1}{\PYZsq{}}\PY{p}{:}\PY{k+kc}{True}\PY{p}{,} \PY{l+s+s1}{\PYZsq{}}\PY{l+s+s1}{CON\PYZus{}NAMES}\PY{l+s+s1}{\PYZsq{}}\PY{p}{:}\PY{k+kc}{True}\PY{p}{,} \PY{l+s+s1}{\PYZsq{}}\PY{l+s+s1}{ASTERISMS}\PY{l+s+s1}{\PYZsq{}}\PY{p}{:}\PY{k+kc}{True}\PY{p}{,} \PY{l+s+s1}{\PYZsq{}}\PY{l+s+s1}{STAR\PYZus{}NAMES}\PY{l+s+s1}{\PYZsq{}}\PY{p}{:}\PY{k+kc}{True}\PY{p}{\PYZcb{}}\PY{p}{)}
\PY{n}{harey}\PY{o}{.}\PY{n}{polar\PYZus{}map}\PY{p}{(}\PY{n}{pole} \PY{o}{=} \PY{l+s+s1}{\PYZsq{}}\PY{l+s+s1}{S}\PY{l+s+s1}{\PYZsq{}}\PY{p}{,} \PY{n}{FOV}\PY{o}{=} \PY{l+m+mi}{110}\PY{p}{,} \PY{n}{mode}\PY{o}{=}\PY{l+s+s1}{\PYZsq{}}\PY{l+s+s1}{azimuth}\PY{l+s+s1}{\PYZsq{}}\PY{p}{)}
\end{Verbatim}
\end{tcolorbox}

    \begin{center}
    \adjustimage{max size={0.9\linewidth}{0.9\paperheight}}{Constellations_memory_demo_files/Constellations_memory_demo_44_0.png}
    \end{center}
    { \hspace*{\fill} \\}
    
    \hypertarget{equatorial-map}{%
\subsection{Equatorial map}\label{equatorial-map}}

The equatorial maps shows the sky around the equator while distorting
the poles. it uses the Gall stereographic projection, which compromises
between preserving the shapes and limit the enlargement of
constellations near the poles.

    \begin{tcolorbox}[breakable, size=fbox, boxrule=1pt, pad at break*=1mm,colback=cellbackground, colframe=cellborder]
\prompt{In}{incolor}{49}{\boxspacing}
\begin{Verbatim}[commandchars=\\\{\}]
\PY{n}{harey}\PY{o}{.}\PY{n}{set\PYZus{}flags}\PY{p}{(}\PY{p}{\PYZob{}}\PY{l+s+s1}{\PYZsq{}}\PY{l+s+s1}{CON\PYZus{}LINES}\PY{l+s+s1}{\PYZsq{}}\PY{p}{:}\PY{k+kc}{True}\PY{p}{,} \PY{l+s+s1}{\PYZsq{}}\PY{l+s+s1}{GRID}\PY{l+s+s1}{\PYZsq{}}\PY{p}{:}\PY{k+kc}{True}\PY{p}{,} \PY{l+s+s1}{\PYZsq{}}\PY{l+s+s1}{CON\PYZus{}NAMES}\PY{l+s+s1}{\PYZsq{}}\PY{p}{:}\PY{k+kc}{True}\PY{p}{,} \PY{l+s+s1}{\PYZsq{}}\PY{l+s+s1}{ASTERISMS}\PY{l+s+s1}{\PYZsq{}}\PY{p}{:}\PY{k+kc}{True}\PY{p}{,} \PY{l+s+s1}{\PYZsq{}}\PY{l+s+s1}{HELPERS}\PY{l+s+s1}{\PYZsq{}}\PY{p}{:}\PY{k+kc}{True}\PY{p}{\PYZcb{}}\PY{p}{)}
\PY{n}{harey}\PY{o}{.}\PY{n}{equatorial\PYZus{}map}\PY{p}{(}\PY{n}{max\PYZus{}dims}\PY{o}{=}\PY{p}{(}\PY{l+m+mi}{10}\PY{p}{,}\PY{l+m+mi}{8}\PY{p}{)}\PY{p}{)}
\end{Verbatim}
\end{tcolorbox}

    \begin{center}
    \adjustimage{max size={0.9\linewidth}{0.9\paperheight}}{Constellations_memory_demo_files/Constellations_memory_demo_46_0.png}
    \end{center}
    { \hspace*{\fill} \\}
    
    \hypertarget{astrolabes-and-planispheres}{%
\subsection{Astrolabes and
planispheres}\label{astrolabes-and-planispheres}}

Astrolabes and planispheres are instruments used to see the stars
visible from a given location during the whole year.

For a complete explanation of how they work, read the
\emph{Astrolabe.ipynb} notebook.

    \begin{tcolorbox}[breakable, size=fbox, boxrule=1pt, pad at break*=1mm,colback=cellbackground, colframe=cellborder]
\prompt{In}{incolor}{50}{\boxspacing}
\begin{Verbatim}[commandchars=\\\{\}]
\PY{n}{harey}\PY{o}{.}\PY{n}{set\PYZus{}flags}\PY{p}{(}\PY{p}{\PYZob{}}\PY{l+s+s1}{\PYZsq{}}\PY{l+s+s1}{CON\PYZus{}LINES}\PY{l+s+s1}{\PYZsq{}}\PY{p}{:}\PY{k+kc}{True}\PY{p}{,} \PY{l+s+s1}{\PYZsq{}}\PY{l+s+s1}{CON\PYZus{}NAMES}\PY{l+s+s1}{\PYZsq{}}\PY{p}{:}\PY{k+kc}{True}\PY{p}{\PYZcb{}}\PY{p}{)}
\PY{c+c1}{\PYZsh{}harey.set\PYZus{}colors(\PYZob{}\PYZsq{}sky\PYZsq{}:\PYZsq{}white\PYZsq{}, \PYZsq{}constellations\PYZsq{}:\PYZsq{}k\PYZsq{}, \PYZsq{}star\PYZsq{}:\PYZsq{}k\PYZsq{}, \PYZsq{}constellation\PYZus{}labels\PYZsq{}:\PYZsq{}red\PYZsq{}\PYZcb{})}
\PY{n}{harey}\PY{o}{.}\PY{n}{create\PYZus{}planisphere}\PY{p}{(}\PY{n}{lat}\PY{o}{=}\PY{l+s+s1}{\PYZsq{}}\PY{l+s+s1}{45 N}\PY{l+s+s1}{\PYZsq{}}\PY{p}{,} \PY{n}{save\PYZus{}folder}\PY{o}{=}\PY{l+s+s1}{\PYZsq{}}\PY{l+s+s1}{planisphere}\PY{l+s+s1}{\PYZsq{}}\PY{p}{,} \PY{n}{FOV}\PY{o}{=}\PY{l+m+mi}{220}\PY{p}{,} \PY{n}{MARK\PYZus{}CENTER}\PY{o}{=}\PY{k+kc}{True}\PY{p}{)}
\end{Verbatim}
\end{tcolorbox}

    \begin{center}
    \adjustimage{max size={0.9\linewidth}{0.9\paperheight}}{Constellations_memory_demo_files/Constellations_memory_demo_48_0.png}
    \end{center}
    { \hspace*{\fill} \\}
    
    \begin{center}
    \adjustimage{max size={0.9\linewidth}{0.9\paperheight}}{Constellations_memory_demo_files/Constellations_memory_demo_48_1.png}
    \end{center}
    { \hspace*{\fill} \\}
    
    \hypertarget{which-constellations-should-you-print}{%
\section{Which constellations should you
print?}\label{which-constellations-should-you-print}}

Not all constellations are the same, both in size and ease of
recognition. Like HARey, I decided to focus on the ones that are bright
and visible from where I live (45° N). Here I made a comparison between
all the constellations, ordering them by visibility.
\textbf{is\_visible()} classify a constellation by where it can be seen.
It considers a ground obstruction/limited visibility zone of 5°.

Constellations are divided in

\begin{itemize}
\tightlist
\item
  visible (during part of the year it is completely above the horizon)
\item
  not visible (always below the horizon during the whole year),
\item
  circumpolar (always visible during the whole year),
\item
  mostly visible (at its highest point, part of the constellation is
  close to the horizon, and light pollution fog and obstacles make it
  difficult to see),
\item
  partly visible (part of the constellation is below the horizon),
\item
  hardly\_visible (part of it is below the horizon, and it never rises
  much from the horizon).
\end{itemize}

    \begin{tcolorbox}[breakable, size=fbox, boxrule=1pt, pad at break*=1mm,colback=cellbackground, colframe=cellborder]
\prompt{In}{incolor}{51}{\boxspacing}
\begin{Verbatim}[commandchars=\\\{\}]
\PY{n}{df} \PY{o}{=} \PY{n}{harey}\PY{o}{.}\PY{n}{stars}\PY{o}{.}\PY{n}{loc}\PY{p}{[}\PY{n}{harey}\PY{o}{.}\PY{n}{stars}\PY{p}{[}\PY{l+s+s1}{\PYZsq{}}\PY{l+s+s1}{constellation}\PY{l+s+s1}{\PYZsq{}}\PY{p}{]} \PY{o}{!=} \PY{l+s+s1}{\PYZsq{}}\PY{l+s+s1}{none}\PY{l+s+s1}{\PYZsq{}}\PY{p}{,} \PY{p}{[}\PY{l+s+s1}{\PYZsq{}}\PY{l+s+s1}{constellation}\PY{l+s+s1}{\PYZsq{}}\PY{p}{,} \PY{l+s+s1}{\PYZsq{}}\PY{l+s+s1}{magnitude}\PY{l+s+s1}{\PYZsq{}}\PY{p}{,} \PY{l+s+s1}{\PYZsq{}}\PY{l+s+s1}{ra}\PY{l+s+s1}{\PYZsq{}}\PY{p}{,} \PY{l+s+s1}{\PYZsq{}}\PY{l+s+s1}{dec}\PY{l+s+s1}{\PYZsq{}}\PY{p}{]}\PY{p}{]}
\PY{n}{df}\PY{p}{[}\PY{l+s+s1}{\PYZsq{}}\PY{l+s+s1}{name}\PY{l+s+s1}{\PYZsq{}}\PY{p}{]} \PY{o}{=} \PY{n}{df}\PY{p}{[}\PY{l+s+s1}{\PYZsq{}}\PY{l+s+s1}{constellation}\PY{l+s+s1}{\PYZsq{}}\PY{p}{]}\PY{o}{.}\PY{n}{map}\PY{p}{(}\PY{n}{harey}\PY{o}{.}\PY{n}{names}\PY{p}{)}
\PY{n}{df}\PY{p}{[}\PY{l+s+s1}{\PYZsq{}}\PY{l+s+s1}{name}\PY{l+s+s1}{\PYZsq{}}\PY{p}{]} \PY{o}{=} \PY{n}{df}\PY{p}{[}\PY{l+s+s1}{\PYZsq{}}\PY{l+s+s1}{name}\PY{l+s+s1}{\PYZsq{}}\PY{p}{]}\PY{o}{.}\PY{n}{str}\PY{o}{.}\PY{n}{replace}\PY{p}{(}\PY{l+s+s1}{\PYZsq{}}\PY{l+s+se}{\PYZbs{}n}\PY{l+s+s1}{\PYZsq{}}\PY{p}{,} \PY{l+s+s1}{\PYZsq{}}\PY{l+s+s1}{ }\PY{l+s+s1}{\PYZsq{}}\PY{p}{)}

\PY{n+nb}{print}\PY{p}{(}\PY{n}{df}\PY{o}{.}\PY{n}{head}\PY{p}{(}\PY{p}{)}\PY{p}{)}

\PY{n}{grouped} \PY{o}{=} \PY{n}{df}\PY{o}{.}\PY{n}{groupby}\PY{p}{(}\PY{p}{[}\PY{l+s+s1}{\PYZsq{}}\PY{l+s+s1}{constellation}\PY{l+s+s1}{\PYZsq{}}\PY{p}{]}\PY{p}{)}
\PY{n}{brightest} \PY{o}{=} \PY{n}{grouped}\PY{p}{[}\PY{l+s+s1}{\PYZsq{}}\PY{l+s+s1}{magnitude}\PY{l+s+s1}{\PYZsq{}}\PY{p}{]}\PY{o}{.}\PY{n}{min}\PY{p}{(}\PY{p}{)}\PY{o}{.}\PY{n}{rename}\PY{p}{(}\PY{l+s+s1}{\PYZsq{}}\PY{l+s+s1}{brightest star}\PY{l+s+s1}{\PYZsq{}}\PY{p}{)}
\PY{n}{mean\PYZus{}luminosity} \PY{o}{=} \PY{n}{grouped}\PY{p}{[}\PY{l+s+s1}{\PYZsq{}}\PY{l+s+s1}{magnitude}\PY{l+s+s1}{\PYZsq{}}\PY{p}{]}\PY{o}{.}\PY{n}{mean}\PY{p}{(}\PY{p}{)}\PY{o}{.}\PY{n}{rename}\PY{p}{(}\PY{l+s+s1}{\PYZsq{}}\PY{l+s+s1}{mean brightness}\PY{l+s+s1}{\PYZsq{}}\PY{p}{)}
\PY{n}{n\PYZus{}stars} \PY{o}{=} \PY{n}{grouped}\PY{o}{.}\PY{n}{size}\PY{p}{(}\PY{p}{)}\PY{o}{.}\PY{n}{rename}\PY{p}{(}\PY{l+s+s1}{\PYZsq{}}\PY{l+s+s1}{n\PYZus{}stars}\PY{l+s+s1}{\PYZsq{}}\PY{p}{)}
\PY{n}{name} \PY{o}{=} \PY{n}{grouped}\PY{p}{[}\PY{l+s+s1}{\PYZsq{}}\PY{l+s+s1}{name}\PY{l+s+s1}{\PYZsq{}}\PY{p}{]}\PY{o}{.}\PY{n}{first}\PY{p}{(}\PY{p}{)}\PY{o}{.}\PY{n}{rename}\PY{p}{(}\PY{l+s+s1}{\PYZsq{}}\PY{l+s+s1}{name}\PY{l+s+s1}{\PYZsq{}}\PY{p}{)}

\PY{n}{visibility} \PY{o}{=} \PY{n}{grouped}\PY{o}{.}\PY{n}{apply}\PY{p}{(}\PY{k}{lambda} \PY{n}{g} \PY{p}{:} \PY{n}{pd}\PY{o}{.}\PY{n}{Series}\PY{p}{(}\PY{n}{harey}\PY{o}{.}\PY{n}{is\PYZus{}visible}\PY{p}{(}\PY{l+s+s1}{\PYZsq{}}\PY{l+s+s1}{45 N}\PY{l+s+s1}{\PYZsq{}}\PY{p}{,} \PY{n}{g}\PY{p}{[}\PY{l+s+s1}{\PYZsq{}}\PY{l+s+s1}{ra}\PY{l+s+s1}{\PYZsq{}}\PY{p}{]}\PY{p}{,} \PY{n}{g}\PY{p}{[}\PY{l+s+s1}{\PYZsq{}}\PY{l+s+s1}{dec}\PY{l+s+s1}{\PYZsq{}}\PY{p}{]}\PY{p}{,} \PY{n}{LVZ} \PY{o}{=} \PY{l+m+mi}{5}\PY{p}{)}\PY{p}{,} \PY{n}{index}\PY{o}{=}\PY{p}{[}\PY{l+s+s1}{\PYZsq{}}\PY{l+s+s1}{Visibility}\PY{l+s+s1}{\PYZsq{}}\PY{p}{,} \PY{l+s+s1}{\PYZsq{}}\PY{l+s+s1}{Best Period}\PY{l+s+s1}{\PYZsq{}}\PY{p}{]}\PY{p}{)}\PY{p}{,} \PY{n}{include\PYZus{}groups}\PY{o}{=}\PY{k+kc}{False}\PY{p}{)}

\PY{n}{constellation\PYZus{}df} \PY{o}{=} \PY{n}{pd}\PY{o}{.}\PY{n}{concat}\PY{p}{(}\PY{p}{[}\PY{n}{name}\PY{p}{,} \PY{n}{brightest}\PY{p}{,} \PY{n}{mean\PYZus{}luminosity}\PY{p}{,} \PY{n}{n\PYZus{}stars}\PY{p}{,} \PY{n}{visibility}\PY{p}{]}\PY{p}{,} \PY{n}{axis}\PY{o}{=}\PY{l+m+mi}{1}\PY{p}{)}
\PY{n}{constellation\PYZus{}df}\PY{o}{.}\PY{n}{head}\PY{p}{(}\PY{l+m+mi}{10}\PY{p}{)}
\end{Verbatim}
\end{tcolorbox}

    \begin{Verbatim}[commandchars=\\\{\}]
     constellation  magnitude        ra        dec        name
hip
677            And       2.07  2.096542  29.090833   Andromeda
746            Cas       2.28  2.292042  59.150222  Cassiopeia
765            Phe       3.88  2.352250 -45.747000     Phoenix
1067           Peg       2.83  3.308958  15.183611     Pegasus
1411           And       9.09  4.418625  37.537861   Andromeda
    \end{Verbatim}

            \begin{tcolorbox}[breakable, size=fbox, boxrule=.5pt, pad at break*=1mm, opacityfill=0]
\prompt{Out}{outcolor}{51}{\boxspacing}
\begin{Verbatim}[commandchars=\\\{\}]
                        name  brightest star  mean brightness  n\_stars  \textbackslash{}
constellation
And                Andromeda            2.07         3.998333       18
Ant                     Pump            4.28         4.280000        1
Aps                      Bee            3.83         4.150000        4
Aql                    Eagle            0.76         3.324545       11
Aqr            Water Carrier            2.90         4.068571       21
Ara                    Altar            2.84         3.506667        6
Ari                      Ram            2.01         4.031000       10
Aur                   Auriga            0.08         2.617500        8
Boo                   Bootes           -0.05         3.493333       15
CMa                  Big dog           -1.44         3.005625       16

                  Visibility           Best Period
constellation
And                  visible       July - December
Ant                  visible        December - May
Aps            never visible                     -
Aql                  visible         May - October
Aqr                  visible       June - November
Ara            never visible                     -
Ari                  visible      August - January
Aur                  visible  September - February
Boo                  visible       February - July
CMa                  visible       October - March
\end{Verbatim}
\end{tcolorbox}
        
    Finally, I tried to order the constellations by how easy they are to be
recognized. I created a simple scoring system based on how many stars
are in the constellation, how bright is the brightest star and the
overall brightness. I decided to weight more the last two scores to
compute the overall score, but this is just my preference.

    \begin{tcolorbox}[breakable, size=fbox, boxrule=1pt, pad at break*=1mm,colback=cellbackground, colframe=cellborder]
\prompt{In}{incolor}{52}{\boxspacing}
\begin{Verbatim}[commandchars=\\\{\}]
\PY{c+c1}{\PYZsh{} Create scoring function to assign scores between 0 and 1}

\PY{k}{def} \PY{n+nf}{score}\PY{p}{(}\PY{n}{x}\PY{p}{)}\PY{p}{:}
    \PY{k}{return} \PY{p}{(}\PY{n}{x}\PY{o}{\PYZhy{}} \PY{n+nb}{min}\PY{p}{(}\PY{n}{x}\PY{p}{)}\PY{p}{)}\PY{o}{/}\PY{p}{(}\PY{n+nb}{max}\PY{p}{(}\PY{n}{x}\PY{p}{)}\PY{o}{\PYZhy{}}\PY{n+nb}{min}\PY{p}{(}\PY{n}{x}\PY{p}{)}\PY{p}{)}

\PY{n}{size\PYZus{}score} \PY{o}{=} \PY{n}{score}\PY{p}{(}\PY{n}{constellation\PYZus{}df}\PY{p}{[}\PY{l+s+s1}{\PYZsq{}}\PY{l+s+s1}{n\PYZus{}stars}\PY{l+s+s1}{\PYZsq{}}\PY{p}{]}\PY{p}{)}
\PY{c+c1}{\PYZsh{} Remap the size score to penalize constellations which are too big or too small}
\PY{n}{size\PYZus{}score} \PY{o}{=} \PY{l+m+mi}{4}\PY{o}{*}\PY{n}{size\PYZus{}score}\PY{o}{*}\PY{p}{(}\PY{l+m+mi}{1}\PY{o}{\PYZhy{}}\PY{n}{size\PYZus{}score}\PY{p}{)}

\PY{c+c1}{\PYZsh{} Mean brightness score, as similar stars are easier to find}
\PY{n}{importance\PYZus{}score} \PY{o}{=} \PY{l+m+mi}{1}\PY{o}{\PYZhy{}}\PY{n}{score}\PY{p}{(}\PY{n}{constellation\PYZus{}df}\PY{p}{[}\PY{l+s+s1}{\PYZsq{}}\PY{l+s+s1}{mean brightness}\PY{l+s+s1}{\PYZsq{}}\PY{p}{]}\PY{p}{)}
\PY{c+c1}{\PYZsh{} Brightest star score}
\PY{n}{brightness\PYZus{}score} \PY{o}{=} \PY{l+m+mi}{1}\PY{o}{\PYZhy{}}\PY{n}{score}\PY{p}{(}\PY{n}{constellation\PYZus{}df}\PY{p}{[}\PY{l+s+s1}{\PYZsq{}}\PY{l+s+s1}{brightest star}\PY{l+s+s1}{\PYZsq{}}\PY{p}{]}\PY{p}{)}

\PY{c+c1}{\PYZsh{} Combine the three scores (I chose to give less importance to the size score)}
\PY{n}{constellation\PYZus{}df}\PY{p}{[}\PY{l+s+s1}{\PYZsq{}}\PY{l+s+s1}{score}\PY{l+s+s1}{\PYZsq{}}\PY{p}{]} \PY{o}{=} \PY{n}{importance\PYZus{}score} \PY{o}{+} \PY{n}{brightness\PYZus{}score} \PY{o}{+} \PY{l+m+mf}{0.8}\PY{o}{*}\PY{n}{size\PYZus{}score}

\PY{c+c1}{\PYZsh{} Keep only the constellations that are visible }
\PY{n}{constellation\PYZus{}df} \PY{o}{=} \PY{n}{constellation\PYZus{}df}\PY{p}{[}\PY{o}{\PYZti{}}\PY{n}{constellation\PYZus{}df}\PY{p}{[}\PY{l+s+s1}{\PYZsq{}}\PY{l+s+s1}{Visibility}\PY{l+s+s1}{\PYZsq{}}\PY{p}{]}\PY{o}{.}\PY{n}{isin}\PY{p}{(}\PY{p}{[}\PY{l+s+s1}{\PYZsq{}}\PY{l+s+s1}{never visible}\PY{l+s+s1}{\PYZsq{}}\PY{p}{,} \PY{l+s+s1}{\PYZsq{}}\PY{l+s+s1}{hardly visible}\PY{l+s+s1}{\PYZsq{}}\PY{p}{]}\PY{p}{)}\PY{p}{]}
\PY{c+c1}{\PYZsh{} Order the constellation by their score}
\PY{n}{constellation\PYZus{}df}\PY{o}{.}\PY{n}{sort\PYZus{}values}\PY{p}{(}\PY{n}{by} \PY{o}{=} \PY{l+s+s1}{\PYZsq{}}\PY{l+s+s1}{score}\PY{l+s+s1}{\PYZsq{}}\PY{p}{,} \PY{n}{ascending}\PY{o}{=}\PY{k+kc}{False}\PY{p}{,} \PY{n}{ignore\PYZus{}index}\PY{o}{=}\PY{k+kc}{False}\PY{p}{)}\PY{o}{.}\PY{n}{head}\PY{p}{(}\PY{l+m+mi}{40}\PY{p}{)}
\end{Verbatim}
\end{tcolorbox}

            \begin{tcolorbox}[breakable, size=fbox, boxrule=.5pt, pad at break*=1mm, opacityfill=0]
\prompt{Out}{outcolor}{52}{\boxspacing}
\begin{Verbatim}[commandchars=\\\{\}]
                         name  brightest star  mean brightness  n\_stars  \textbackslash{}
constellation
CMa                   Big dog           -1.44         3.005625       16
Pup                     Stern           -0.62         3.291000       10
Aur                    Auriga            0.08         2.617500        8
Boo                    Bootes           -0.05         3.493333       15
Sco                  Scorpion            1.06         2.910500       20
Aql                     Eagle            0.76         3.324545       11
Cyg                      Swan            1.25         3.386000       15
Lyr                      Lyre            0.03         3.714000       10
Vir                    Maiden            0.98         3.562667       15
Gem                     testa            1.16         3.345556       18
Leo                      Lion            1.36         3.361250       16
UMa                  Big Bear            1.76         2.991579       19
CMi                Little dog            0.40         1.645000        2
Tau                      Bull            0.87         3.571000       20
Sgr                    Archer            1.79         3.313333       18
Ori                     Orion            0.18         3.410435       23
Hya                     Hydra            1.99         3.585333       15
Dra                    Dragon            2.24         3.522667       15
Gru                     Crane            1.73         3.444444        9
Peg                   Pegasus            2.38         3.645882       17
Oph            Serpent bearer            2.08         3.742222       18
Cet                     Whale            2.04         3.866154       13
Cas                Cassiopeia            2.15         2.771667        6
Cep                   Cepheus            2.45         3.648182       11
Lep                      Hare            2.58         3.618182       11
And                 Andromeda            2.07         3.998333       18
PsA             Southern Fish            1.17         4.021250        8
Ari                       Ram            2.01         4.031000       10
UMi              Little  Bear            1.97         3.548571        7
Per                   Preseus            1.79         3.649130       23
Ser                   Serpent            2.63         3.789500       20
Cap                 Capricorn            2.85         4.010000       13
Crv                     Raven            2.58         3.530000        7
Lib                    Scales            2.61         3.296667        6
Col                      Dove            2.65         3.738571        7
CrB            Northern Crown            2.22         3.760000        6
Cen                   Centaur           -0.01         3.330345       29
Aqr             Water Carrier            2.90         4.068571       21
Her                  Hercules            2.78         3.753913       23
Lup                      Wolf            2.30         3.856667       24

                   Visibility           Best Period     score
constellation
CMa                   visible       October - March  2.392003
Pup            partly visible       October - March  2.086273
Aur                   visible  September - February  2.074486
Boo                   visible       February - July  2.042749
Sco            mostly visible        March - August  1.937943
Aql                   visible         May - October  1.901712
Cyg                   visible         May - October  1.874205
Lyr                   visible     April - September  1.865028
Vir                   visible        January - June  1.864850
Gem                   visible       October - March  1.862896
Leo                   visible        December - May  1.860364
UMa                   visible        December - May  1.844017
CMi                   visible       October - March  1.815064
Tau                   visible  September - February  1.777285
Sgr            mostly visible     April - September  1.775529
Ori                   visible  September - February  1.770069
Hya                   visible        December - May  1.703428
Dra               circumpolar                Always  1.683092
Gru            partly visible       June - November  1.636852
Peg                   visible       June - November  1.609887
Oph                   visible        March - August  1.607807
Cet                   visible      August - January  1.598738
Cas               circumpolar                Always  1.582088
Cep               circumpolar                Always  1.549510
Lep                   visible     September - March  1.538192
And                   visible       July - December  1.535746
PsA                   visible       June - November  1.503932
Ari                   visible      August - January  1.470255
UMi               circumpolar                Always  1.455835
Per                   visible      August - January  1.454546
Ser                   visible        March - August  1.444559
Cap                   visible         May - October  1.433169
Crv                   visible        January - June  1.367613
Lib                   visible       February - July  1.360674
Col            mostly visible       October - March  1.296943
CrB                   visible        March - August  1.287348
Cen            partly visible        January - June  1.283449
Aqr                   visible       June - November  1.278057
Her                   visible        March - August  1.272596
Lup            partly visible       February - July  1.247301
\end{Verbatim}
\end{tcolorbox}
        

    % Add a bibliography block to the postdoc
    
    
    
\end{document}
